% !TEX root = ../thesis_main.tex

% Kapitel 6 — Budget: 10 Seiten — Schreiben: 19.10.–13.11.; 30.11.–11.12. auf den gefahrenen Stand bringen (Generalprobe, Pilot, Sitzungen).
% Bögen, Prüflisten und Leitfaden in den Anhang.
%
% ===========================================================================
% GESAMTKONZEPT KAPITEL 6
% Notizen am 05.11. neu gefasst. Jede Angabe gegen das Werkzeug-Repo (codebase-graph-mcp) und den Studienrechner (studie\) geprüft; was nicht belegt ist, trägt KLÄREN. ===========================================================================
%
% STAND DER QUELLEN
% - Werkzeug-Repo: letzter Studien-Commit 628b04f (02.11., Ablauf vor der ersten Sitzung neu gegliedert; Bögen: Teil A ans Ende, TLX-Anker; Interviewfrage 4; Protokoll 29.10.). Nicht committet: eval/studie/ablauf-vorlage.md (Screenshot zum Stummschalten in RustDesk geht vor der Sitzung an den Probanden) und sitzung-anlegen.ps1 (kopiert ihn).
% - eval/studie/einwilligung.md ist im Repo unverändert und trägt noch alle Platzhalter. Die verschickte Fassung liegt als PDF in den Rückläufen unter studie\einwilligung\ (02.11., 05.11., 05.11.) und steht seit dem 07.11. im Wortlaut in 6.5 WORTLAUT; gefüllte Platzhalter siehe SCHREIBGERÜST unten.
% - Sitzungsstand (studie\probanden\<Kennung>-sitzung\01-mitschrieb.md): Pilot PG 22.10., P03 03.11., P04 05.11. P01 war nur Durchspielen (21./22.10.), P02 hat einen Sitzungsordner (23.10.), ist nicht gelaufen. Auswertungspipeline unter studie\auswertung\ (29.10.), Ergebnisse für PG, P03, P04 liegen dort.
% - Bei Widerspruch gilt: Mitschrieb (gefahren) vor Repo-Dateien (Bögen, Ablauf, Skripte) vor Docs (Plan, Mess). Kapitel 6 beschreibt den gefahrenen Stand, 7.1 die Abweichungen.
%
% QUELLENKÜRZEL (Werkzeug-Repo, wenn nicht anders gesagt)
%   Plan    = docs/studie-plan.md (Stand 11.10., Nachträge bis 02.11.)
%   Mess    = docs/studie-messung-entwurf-2026-10-08.md (Herleitung; Wertung dort überholt)
%   Leit    = docs/studie-leitfaden-sitzungsleitung.md (Stand 21.10.)
%   Ablauf  = eval/studie/ablauf-vorlage.md, ablauf-aufgabe-vorlage.md (Stand 02.11.)
%   Zut     = eval/studie/kopien-zuteilung.ps1
%   T74, TE = eval/tasks/teamable-frontend-74|E/ (task.md, vorher-/nachher-bogen.md, schluessel.md, erklaeren.md, pruefliste.md, prompts.md, reference-context.json)
%   TR      = eval/tasks/README.md (Auswahlregeln, Stand 10.09.)
%   CLauf   = eval/results-studie-c-opus-20261009/README.md (C-Läufe 09.10., Rechnung 13.10.)
%   Prot    = docs/arbeitsprotokoll.md
%   Prüf    = docs/studie-pruefung-2026-10-11.md
%   Gl      = docs/notizen/Zeitplan und Gliederung 2026-10-19.md
%   Tr19    = docs/notizen/Notizen Meeting 2026-10-19.md
%   N16     = docs/notizen/Notizen Meeting 2026-09-16.md (Mail 16.09., Antwort 13.10.)
%   Zb      = docs/notizen/Zu besprechen.md (24.09., Stufe 3)
%   Ex      = Expose_v4/sections/02_forschungsfragen.tex, 03_methodik.tex
%   Ausw    = studie\auswertung\README.md (Transkription, 29.10.)
%   Mit     = studie\probanden\<Kennung>-sitzung\01-mitschrieb.md
%
% LEITGEDANKE (roter Faden, den jeder Abschnitt trägt)
% Die Studie prüft nicht mehr, ob kuratierter Kontext den Agenten besser macht, sondern ob der Mensch, der den Kontext selbst zusammenstellt, die Änderung des Agenten danach besser versteht. Der Weg dahin steht in den Benchmark-Messungen: AP8 04.09. (mit Graph weniger Suchen, gefunden wird dasselbe) und die Rechnung vom 13.10. (der automatische Kontext deckt den Referenzkontext fast vollständig). Das Design folgt daraus Punkt für Punkt: C gegen D statt ohne Werkzeug gegen D (6.2), zwei Aufgaben mit bekanntem Spielraum (6.3), Fragen vorher und nachher statt nur nachher (6.4), Messwiederholung mit acht Personen (6.5), Richtung statt Nachweis (6.6).
%
% KAPITELBOGEN UND ABGRENZUNG
% - 1.5 nennt FF1 bis FF6 im Wortlaut des Exposés, unverändert (Antwort des Betreuers 13.10.: nicht umdrehen, Wechsel begründen, Entwicklung im Fazit; N16).
% - 5.6 zieht Bilanz der Benchmark-Messungen und begründet den Wechsel der Hauptfrage aus den Zahlen. 6.1 setzt dort an, wiederholt die Zahlen nicht (nur Verweis), ordnet die zwei Fragen der Studie den FF zu und sagt in zwei Sätzen, was vom Exposé-Design entfällt.
% - Abweichungen vom Exposé bekommen kein Sammelkapitel; jede steht in zwei Sätzen dort, wo der Leser sie zuerst trifft, mit der Zahl, die sie ausgelöst hat (Gl): Backend zu Frontend (6.3), Stufe 3 gestrichen (6.1), C statt ohne Werkzeug (6.2), keine Zweitableitung durch den Betreuer (6.3), 100 statt 60 Minuten (6.5).
% - 7.1 berichtet, was bei der Durchführung vom Plan abwich (Dauer, Störungen, Abbrüche). 8.2 sammelt die Grenzen ein, die Kapitel 6 je Abschnitt benennt (Liste unten).
% - Anhang: Aufgabenblätter (app:aufgaben); Vorfragebogen, Vorher- und Nachher-Bögen, Schlüssel, Aussagenlisten, Prüflisten, SUS, Interviewfragen, Leitfaden, Ablauf, Einwilligung (app:frageboegen). Im Text nur Aufbau, Tabellen und je ein Beispielitem.
%
% ENTSCHEIDUNG FF-NUMMERN (vor 6.1 festlegen)
% Der Plan nummeriert "FF6 explorativ -> Frage 2, Frage 1 nimmt FF4 auf, FF2/FF3 = Wirkung auf den Agenten, Aufwand = FF5". Das passt zu keiner Fassung des Exposés v4 (Stand 12.08. und 27.10. geprüft, Nummern gleich): FF1 Gestaltung mit E1 bis E3, FF2 Nähe der Entwicklerauswahl zur optimierten Auswahl (D gegen D'), FF3 Kontextgüte (C gegen D), FF4 Lösungsqualität, Schrittzahl, Tokens, FF5 strukturiert gegen flach (B gegen C), FF6 Interaktionsaufwand und Bewertung. 1.5 übernimmt v4. Vorschlag für 6.1 (v4-Nummern):
%   Frage 1 (Eingriffe; wird der Kontext besser oder schlechter) = FF3, mit FF2 als Bezug auf den Referenzkontext (D' wurde nicht gebaut; der Spielraum ist per Rechnung 13.10. bekannt und je Aufgabe benannt, 6.3).
%   Frage 2 (Verständnis) = neu. Herkunft: Nebennutzen in FF6 ("wer den Kontext selbst zusammenstellt, kennt die Codebase hinterher womöglich besser", Ex) und das Transparenzproblem (1.4). E3 (FF1) wird mit erhoben.
%   Nebenfrage Aufwand und Bewertung = FF6 (Kuratierzeit, TLX, SUS, Interview).
%   Nebenfrage Ergebnis des Agenten = FF4, nur C gegen D je Aufgabe über die Prüfliste.
%   FF5 entfällt mit Stufe 3; dazu nur Verweis auf AP8 (5.3; C war dort nicht verschachtelt, Zb Punkt 2) und auf die Rückgabeform (4.5).
% Die Plan-Nummern aus den Repo-Docs nicht übernehmen.
%
% BEGRIFFE (Kapitel 4 sagt "Nutzer", Plan und Bögen "Proband", Exposé "Entwickler")
% Festlegen, Vorschlag: "Teilnehmende" für die Gruppe, "Proband" für die Person in der Sitzung, "Sitzungsleitung" für den Autor in dieser Rolle, "Agent" für Claude Code, "die zwei Entwickler" für die Vergleichsfälle. Bedingungen immer "C" und "D", beim ersten Auftreten in 6.2 erklärt. Oberfläche = "Oberfläche des stufenweisen Modus" wie in 4.5, in der Studie die "Studienseite".
%
% GRENZEN, die Kapitel 6 je Abschnitt benennt und 8.2 einsammelt
% 6.1 Stufe 3 entfällt, FF4 und FF5 bleiben ohne eigene Messung.
% 6.3 Referenzkontext aus einer Hand, keine Zweitableitung, E ohne PR und Text in dritter Fassung, sealedIn offen, Deckeneffekt bei #74, Eigenschaft 2 des Exposés ungeprüft.
% 6.4 Vortest als Suchauftrag; Wertung des Erklärens nicht verblindet; Items teils aus der Zusammenfassung lösbar (Schicht "geliefert"); Kontrollitem E5 bei Opus geliefert; laute Antwort über den Fernzugriff (absolute Kalibrierung nicht belastbar); Sehen und Wählen in D untrennbar; Vertrauensskala des Exposés nicht erhoben.
% 6.5 acht Personen im zweiten Sprint, Pilot mit altem Ablauf, Bögen über der Zeitgrenze.
% 6.6 kein Test, Richtung je Person.

% ---------------------------------------------------------------------------
% SCHREIBGERÜST (07.11.)
% Unter jedem Abschnitt stehen seit dem 07.11. die Blöcke GERÜST (Absatzplan, Tabellen, Beispielitems) und WORTLAUT (Aufgabentexte, Prompts, Bögen mit Lösung, Schlüsselbefunde, Aussagenlisten, Prüflisten, Ablauf mit Ansagen, Einwilligung in der verschickten Fassung, Zuteilung, Mitschriebe). Alles aus dem Werkzeug-Repo und dem Studienordner, was Kapitel 6 und die Anhänge app:aufgaben und app:frageboegen brauchen, steht damit in dieser Datei; das Kapitel lässt sich ohne Repo und Studienrechner schreiben. Nicht enthalten: was nach dem 07.11. gefahren wird (P05 bis P10, P01, P02); das kommt beim Nachziehen 30.11. bis 11.12. aus den Mitschrieben (Vorlage des Mitschriebs in 6.5 WORTLAUT).
% STAND 07.11. der bisher offenen Punkte:
% - Einwilligung: verschickte Fassung liegt vor (PDF im Rücklauf, Wortlaut in 6.5 WORTLAUT). Gegenüber dem Repo-Stand gefüllt: Fakultät Informatik, HTWG Konstanz, Betreuer Prof. Dr. Rainer Mueller, Mail des Autors, Datenverantwortlicher der Autor mit Anschrift, Transkription lokal, Call über Discord (Discord Netherlands BV, Standardvertragsklauseln und EU-US Data Privacy Framework), Beschwerde beim Landesbeauftragten für den Datenschutz Baden-Württemberg, Löschung Aufnahmen spätestens 16.02.2027, übrige Daten spätestens 30.09.2027, Rücktrittsfrist 11.12.2026, letzter Sitzungstermin 27.11.2026; Zeile zur Datenschutzbeauftragten der Hochschule gestrichen; "acht Personen aus einem Teamprojekt" zu "größtenteils aus demselben Teamprojekt"; bei Antwort per Mail soll die Person dazusagen, ob wörtliche Zitate erlaubt sind. Drei Rückläufe: 02.11. unterschriebenes PDF, alle fünf Punkte angekreuzt; 05.11. unterschriebenes PDF; 05.11. Mail "einverstanden" ohne Angabe zu Zitaten. eval/studie/einwilligung.md im Repo bleibt die Platzhalterfassung.
% - Markierung der nötigen Stellen durch die zwei Entwickler: nicht geschehen. Mitschrieb PG 22.10. ohne Eintrag, keine Markierungsdatei in den Sitzungsordnern PG, P01, P02; P01 und P02 sind nicht gelaufen (P01 nur Durchspielen 21./22.10. ohne Probanden, P02 Ordner vom 23.10.). Im Text so geschrieben (6.3).
% - Auslese-Skript für die Bögen: gibt es nicht. studie\auswertung\ enthält nur aufnahme-auswerten.py und sprecher-zuordnen.py. Die Bögen liegen als Markdown im Sitzungsordner (05- bis 09-) und werden von Hand ausgelesen; so in 6.6 schreiben.
% - Zweite Gruppe: am 29.10. entschieden, nur dieses eine Teamprojekt, mit dem Betreuer abgesprochen (Notizen der Vorstellung). Die Zeile "Teamprojekt 2, sieben Personen, 16. bis 27.11." im Zeitplan des Plans ist damit hinfällig; n bleibt acht plus zwei Entwickler.
% - Sitzungstermine: Werktage 30.10. bis 27.11., Blöcke von zwei Stunden (100 min plus 15 min Vorlauf), eine Person je Block, zwischen zwei Sitzungen eine Stunde frei; möglichst bis 13.11., danach Reserve.
% ---------------------------------------------------------------------------

\chapter{Nutzerstudie}
\label{chap:nutzerstudie}

\section{Fragen der Studie und ihre Herkunft aus den Benchmark-Messungen}% 1,5 S.
\label{sec:studienfragen}
% ZIEL: Der Leser kennt die zwei Fragen der Studie, ihre Zuordnung zu den FF und den Grund, warum die Studie Verständnis misst und nicht die Leistung des Agenten. Keine Messzahlen aus Kapitel 5 wiederholen, nur verweisen. 6.1 schreiben wir gemeinsam.
%
% INHALT (geprüft)
% - Frage 1: Was macht der Proband mit dem Kontext, und wird der Kontext dadurch besser oder schlechter? Kopfzahl: Anteil der Referenzstellen im übergebenen Kontext, kuratiert minus automatisch, je Person (Plan E1).
% - Frage 2 (Hauptfrage): Kann er danach besser erklären und beurteilen, was der Agent getan hat? Kopfzahl: Zuwachs vorher zu nachher auf nicht gelieferten Items, D minus C je Person (Plan E1, E3).
% - Verbindung beider Fragen: je Item und Person, ob der Knoten, an dem das Item hängt, gezeigt, angesehen, behalten oder abgewählt wurde, und ob das Item richtig ist (Plan).
% - Mit erhoben: Kalibrierung (sicher und falsch). Nur beschreibend: ob der Proband weiß, worauf der Agent sich stützt; in D per Konstruktion bekannt, daher keine Punktzahl (Plan, Mess Befund "nahezu per Konstruktion wahr").
% - Nebenfragen: Aufwand (Kuratierzeit, TLX, SUS, Interview) und Ergebnis des Agenten (Prüfliste je Aufgabe, C gegen D) (Plan).
% - E3 aus 1.5 (Begründungen werden beim Kuratieren genutzt) wird über den Export der Oberfläche (angesehene Knoten, Eingriffe), das laute Denken und Interviewfrage 3 geprüft (01_einleitung.tex, Ex).
%
% HERKUNFT, in dieser Reihenfolge erzählen
% 1. Exposé (Ex 03_methodik): vier Stufen; Bedingungen A, B, C, D, D'; Stufe 4 = Studie C gegen D als Messwiederholung mit acht Studierenden, 60 Minuten, TLX, Vertrauen, SUS. Hauptgewicht lag auf Stufe 3 (Agent löst Aufgaben automatisch unter A, B, C, D').
% 2. AP8 04.09. (5.3): ohne Kontext gegen automatischen Kontext ist gemessen (weniger Suchen, gefunden wird dasselbe). Folge 02.10.: C gegen D statt ohne Werkzeug gegen D (Plan E2, Tr19 Punkt 3).
% 3. Stufe 3 entfällt: Tokenbudget und erwarteter Nullbefund auf der kleinen Codebasis (Zb 24.09.: 36 Läufe, 25 bis 110 USD; N16 Mail 16.09.; Tr19). Am 19.10. aus Zeitgründen nicht besprochen, danach eigene Entscheidung: entfällt (Tr19 Mitschrift, Plan Offen). Damit nicht belegt: kuratierter Kontext hilft dem Agenten mehr als automatischer; FF4 nur noch C gegen D je Person über die Prüfliste, FF5 ohne neue Messung (Zb Punkt 1).
% 4. Rechnung 13.10. (CLauf, 5.3 oder 5.6): Deckung des automatischen Kontexts bei #74 in allen drei Anfragen 7 von 7, bei E 9, 8, 8 von 9. Kuratieren kann den Kontext kaum verbessern; Frage 1 bleibt beschreibend, Hauptfrage wird das Verständnis (Plan, Tr19 Punkt 2).
% 5. Betreuer: Mail 16.09. mit dem Vorschlag, die Hauptfrage umzudrehen (Grund: Agent findet mit und ohne Kontext dieselben Stellen; das Unterscheidende ist der Mensch, der Schritt für Schritt sieht und wählt). Antwort 13.10.: Richtung passt, FF des Exposés stehen lassen, Wechsel begründen, Entwicklung im Fazit beschreiben (N16). 19.10.: Studiendesign durchgesprochen, passt so (Tr19 Mitschrift, Plan).
% 6. Messung 08.10.: Fassung vom 02.10. (zwei vorbereitete Diffs, fünf Fragen nach dem Lauf) verworfen, weil die Items aus Aufgabentext, Diff und Zusammenfassung des Agenten lösbar waren, in C und D gleich; daher Vorher-Nachher (Mess; Einzelheiten 6.4).
%
% VORAB FESTLEGEN, was gegen die Annahme spricht (Ex, letzter Absatz "Zuordnung")
% - Frage 2: "Kein erkennbarer Unterschied" ist ein vollwertiges Ergebnis mit der Lesart, dass Kuratieren kein Verständnis über eine gute Zusammenfassung des Agenten hinaus bringt (Plan, Mess).
% - Frage 1: bei #74 ist nur null oder eine Verschlechterung möglich; bei E höchstens die eine fehlende Datei (6.3). Berichtet wird deshalb je Aufgabe getrennt.
% - Exposé-Kriterium (Ausschöpfungsgrad nahe null und kein Vorteil bei den Akzeptanzkriterien) bleibt für FF2 bis FF4 benannt, trägt aber nicht mehr die Arbeit.
%
% VERWEISE: 1.4 (Transparenzproblem), 1.5 (FF, E1 bis E3), 4.5 (Betriebsmodi, Export), 5.3 (AP8, Läufe 05.10. und 09.10.), 5.6 (Bilanz, Begründung), 8.1 (Antworten je FF).
% Text am 05.11. geschrieben, FF-Zuordnung nach dem Vorschlag im Kapitelkopf (Beschluss 05.11.). Die Zahlen der Rechnung vom 13.10. stehen hier in einem Satz, ausführlich in 6.3. Begriffe: Proband = Nutzer aus Kapitel 4 in der Sitzung, Sitzungsleitung = Autor.
% Neufassung 06.11.: Daten, Mails, Betreuertreffen und die Stufenliste des Exposés aus der Prosa genommen (Chronik gehört laut Betreuer ins Fazit, die Notizen oben behalten sie). Jede Designentscheidung steht jetzt mit ihrem Befund, nicht mit ihrem Datum. Das Exposé wird nur noch einmal indirekt genannt ("ursprünglicher Plan"). Gestrichen: Satz zum Exposé-Kriterium (Ausschöpfungsgrad nahe null); gehört zu 8.1, wo FF2 bis FF4 beantwortet werden. Ergänzt 06.11. (aus dem Gutachten): Tabelle FF-Zuordnung statt des dichten Absatzes, Wirkannahme im Hauptfrage-Absatz, Definition von Verständnis im Vorher-Nachher-Absatz, Sehen und Wählen nicht trennbar bei "was als Ergebnis gilt".

% GERÜST 6.1 (geschrieben 05./06.11., acht Absätze, steht): 1 zwei Fragen und Begriffe; 2 Herkunft aus den FF, Tabelle tab:ff-zuordnung; 3 warum Verständnis (Deckung in einem Satz, Wirkannahme); 4 warum C gegen D; 5 warum keine automatisierten Läufe (Stufe 3); 6 Definition Verständnis, vorher und nachher; 7 was als Ergebnis gilt; 8 Vorschau auf 6.2 bis 6.6. Der Wortlaut von FF1 bis FF6 steht in chapters/01_einleitung.tex (Thesis-Repo), nichts davon liegt nur im Werkzeug-Repo.
% WORTLAUT der zwei Fragen im Plan (E1), falls die Prosa noch einmal daran geprüft werden soll: "Frage 1: Was macht der Proband mit dem Kontext, den das Werkzeug ihm vorschlägt, und wird der Kontext dadurch besser oder schlechter? Frage 2 (Hauptfrage): Kann er danach besser erklären und beurteilen, was der Agent getan hat?" Kopfzahlen: Frage 1 Anteil der Referenzstellen im übergebenen Kontext, kuratiert minus automatisch, je Person; Frage 2 Zuwachs vorher zu nachher auf nicht gelieferten Items, D minus C je Person.

Die Studie stellt zwei Fragen. Die erste fragt nach dem Kontext: Was macht ein Proband mit dem Kontext, den das Werkzeug ihm vorschlägt, und wird der Kontext dadurch besser oder schlechter? Die zweite fragt nach dem Verständnis: Kann er danach besser erklären und beurteilen, was der Agent getan hat? Die zweite Frage ist die Hauptfrage der Studie. Proband heißt in diesem Kapitel die Person, die in einer Sitzung die Rolle des Nutzers aus Kapitel~\ref{chap:werkzeug} einnimmt. Sitzungsleitung ist der Autor.

Beide Fragen gehen aus den Forschungsfragen in Abschnitt~\ref{sec:forschungsfragen} hervor, decken sich aber nicht mit ihnen. Tabelle~\ref{tab:ff-zuordnung} ordnet jede Forschungsfrage dem zu, was die Studie daraus macht. Die erste Frage nimmt FF3 auf, die Veränderung der Kontextgüte gegenüber dem automatischen Abruf. Die zweite Frage steht so nicht unter den Forschungsfragen. Sie greift einen Nebengedanken aus FF6 auf: Wer den Kontext selbst zusammenstellt, setzt sich dabei mit der Codebasis auseinander und kennt sie hinterher womöglich besser. Aus diesem Nebengedanken wird die Hauptfrage.

\begin{table}[htbp]
  \centering
  \caption{Forschungsfragen aus Abschnitt~\ref{sec:forschungsfragen} und ihre Behandlung in der Studie.}
  \label{tab:ff-zuordnung}
  \small
  \begin{tabularx}{\textwidth}{l>{\raggedright\arraybackslash}X}
    \toprule
    Forschungsfrage & In der Studie \\
    \midrule
    FF1, E3 & Mit erhoben: ob die Probanden die offengelegten Begründungen beim Kuratieren heranziehen. E1 und E2 belegt Kapitel~\ref{chap:werkzeug}. \\
    FF2 & Liefert den Maßstab, den Referenzkontext je Aufgabe. Die optimierte Auswahl wurde nicht gebaut; an ihre Stelle tritt eine Rechnung, die den Spielraum des Kuratierens je Aufgabe bestimmt (Abschnitt~\ref{sec:aufgaben}). \\
    FF3 & Frage 1 der Studie: Veränderung der Kontextgüte gegenüber dem automatischen Abruf, beschreibend und je Aufgabe getrennt. \\
    FF4 & Nebenfrage: Ergebnis des Agenten je Person und Aufgabe aus einem einzelnen Lauf, über eine Prüfliste (Abschnitt~\ref{sec:messung}). \\
    FF5 & Nicht untersucht; bleibt beim Befund aus Abschnitt~\ref{sec:agentenverhalten}. \\
    FF6 & Nebenfragen Aufwand und Bewertung (Kuratierzeit, Beanspruchung, Gebrauchstauglichkeit, Interview). Der dort genannte Nebennutzen wird Frage 2, die Hauptfrage. \\
    \bottomrule
  \end{tabularx}
\end{table}

Dass die Hauptfrage auf das Verständnis zielt und nicht auf den Kontext, folgt aus den Benchmark-Messungen (Abschnitt~\ref{sec:ueberleitung-studie}). Für beide gewerteten Aufgaben wurde vor der Studie an den Anfragen des Agenten gerechnet, welchen Teil des Referenzkontexts der automatische Abruf schon liefert. Bei der ersten Aufgabe deckte er in allen drei geprüften Anfragen jede Stelle des Referenzkontexts, bei der zweiten acht oder neun von neun Stellen (Abschnitt~\ref{sec:aufgaben}). Kuratieren kann den Kontext auf diesen Aufgaben also kaum verbessern, bei der ersten Aufgabe nur gleich lassen oder verschlechtern. Die erste Frage wird deshalb beschreibend beantwortet. Was das Werkzeug von einem automatischen Abruf unterscheidet, ist nicht der Kontext, der am Ende beim Agenten ankommt, sondern der Mensch, der ihn Schritt für Schritt sieht und auswählt. Die Annahme dahinter: Wer den Kontext kuratiert, hat sich vor dem Lauf mit den Stellen beschäftigt, an denen der Agent danach arbeitet, und kann dessen Änderung deshalb einordnen. Wer im Terminal zusieht, trifft diese Stellen erst im Ergebnis. Ob dieser Mensch danach die Änderung des Agenten besser versteht, können die Benchmark-Messungen nicht beantworten. Das ist die Frage, für die es Probanden braucht.

Aus demselben Grund vergleicht die Studie zwei Bedingungen, die sich nur in der Beteiligung des Probanden unterscheiden. In beiden arbeitet der Agent mit dem Kontext aus dem Graphen. In C erhält er ihn automatisch, in D hat ein Mensch ihn vorher gesehen und verändert (Abschnitt~\ref{sec:bedingungen}). Ein Vergleich ohne Werkzeug gegen D läge näher, hätte aber einen Unterschied noch einmal gemessen, der aus Abschnitt~\ref{sec:agentenverhalten} schon bekannt ist: Mit dem Kontextblock sucht der Agent weniger, findet am Ende jedoch dasselbe. Dieser bekannte Unterschied hätte sich mit der Beteiligung des Probanden vermischt.

Auf automatisierte, mehrfach wiederholte Läufe des Agenten unter verschiedenen Arten von Kontext, wie sie der ursprüngliche Plan als eigene Stufe vorsah, verzichtet die Arbeit. Sie hätten je Aufgabe und Bedingung mehrere vollständige Agentenläufe verlangt, und nach dem Befund aus Abschnitt~\ref{sec:agentenverhalten} war auf der kleinen Codebasis der Studie kein messbarer Vorteil des kuratierten Kontexts für den Agenten zu erwarten. Damit bleibt unbelegt, ob kuratierter Kontext dem Agenten mehr hilft als automatischer. FF4 und FF5 bleiben deshalb ohne eigene Messung (Tabelle~\ref{tab:ff-zuordnung}).

Verständnis meint dabei nicht die Codebasis im Ganzen, sondern das, was die Bögen und das Erklären abfragen: welche Stellen für die Aufgabe zusammenhängen, was der Agent geändert hat und warum (Abschnitt~\ref{sec:messung}). Gemessen wird es zweimal, vor dem Lauf des Agenten und danach, und gewertet wird der Zuwachs. Fragen allein nach dem Lauf reichen nicht: Eine Prüfung an vier Agentenläufen zeigte, dass sich solche Fragen aus dem Aufgabentext, dem Diff und der Zusammenfassung des Agenten beantworten ließen, in C und D gleichermaßen. Abschnitt~\ref{sec:messung} begründet das im Einzelnen.

Was als Ergebnis gilt, ist vorab festgelegt. Für die Hauptfrage ist ein Vorsprung von D gegenüber C beim Zuwachs an richtig beantworteten Fragen die erwartete Richtung. Ein Vorsprung ist dabei dem Verfahren als Ganzem zuzurechnen: Wer kuratiert, hat mehr Code gesehen und zugleich gewählt, und die Studie trennt beides nicht (Abschnitt~\ref{sec:bedingungen}). Zeigt sich kein erkennbarer Unterschied, ist auch das ein vollwertiges Ergebnis mit einer klaren Lesart: Das Kuratieren bringt dann kein Verständnis über eine gute Zusammenfassung des Agenten hinaus. Für die erste Frage ist bei der ersten Aufgabe nur null oder eine Verschlechterung möglich, bei der zweiten höchstens das Nachholen einer fehlenden Datei. Sie wird deshalb je Aufgabe getrennt berichtet. Bei acht Personen berichtet die Studie Richtungen je Person und keinen Nachweis (Abschnitt~\ref{sec:auswertung}).

Die folgenden Abschnitte beschreiben die beiden Bedingungen und die Einstellung von Claude Code dafür (Abschnitt~\ref{sec:bedingungen}), die Aufgaben und ihren Referenzkontext (Abschnitt~\ref{sec:aufgaben}), die Messung (Abschnitt~\ref{sec:messung}), die Teilnehmenden und den Ablauf einer Sitzung (Abschnitt~\ref{sec:teilnehmende}) und die Auswertung (Abschnitt~\ref{sec:auswertung}).

\section{Bedingungen C und D}% 1 S.
\label{sec:bedingungen}
% ZIEL: Der Leser weiß, was in C und D gleich ist und was sich allein unterscheidet, wie die Zeit je Bedingung läuft und wie Claude Code dafür eingestellt ist. Text am 06.11. geschrieben, danach von drei auf etwa eine Seite gekürzt: vier Absätze, Einstellungen und Wartezustand (Hintergrund, Abbruch, Leerlauffrist; am 04.11. aus 4.5 hierher verschoben) in zwei Fußnoten.
%
% INHALT (geprüft)
% - C: Claude Code mit get_context; der Agent bekommt den automatischen Kontext (30 Knoten, Seeds zuerst, hopDepth 2; 4.4, 4.5), der Proband sieht im Terminal zu. D: Claude Code mit get_context_stepwise; der Proband kuratiert in der Studienseite der Oberfläche, bestätigt, danach arbeitet der Agent (Plan E2).
% - In beiden Bedingungen Graph-Kontext, Unterschied nur die Beteiligung des Probanden. Am 02.10. umgestellt von "ohne Werkzeug gegen D" (Plan E2; Begründung 6.1).
% - Verweis aus 4.2 hier einlösen: Im stufenweisen Modus lädt der Proband die Nachbarn selbst (Stufe 3 der Oberfläche), ein D-Kontext ohne Eingriff enthält nur die Seeds. Deshalb ist der Vergleichswert für Frage 1 nicht der Startzustand der Oberfläche, sondern get_context neu gerechnet mit task und userQuery aus dem Export derselben Sitzung (Plan E3; 6.4).
% - Sehen und Verändern sind nicht trennbar: Wer in D kuratiert, hat mehr Code gesehen. Dritte Bedingung "D nur lesend" verworfen (Mess Option D: 2 bis 3 Personen je Zelle, Lesemodus und dritte Aufgabe nötig). Kuratierzeit wird als Kosten berichtet (Plan).
% - Zeit je Aufgabe (Ablauf, Stand 02.11.): in beiden Bedingungen erst Aufgabenblatt 2 min, dann Vorher-Bogen bis Minute 4, dann drückt der Proband Enter. In C arbeitet der Agent sofort. In D kuratiert der Proband bis zur Bestätigung, höchstens 15 min (Ansage bei 13, Bitte bei 15; sitzung-anlegen.ps1, laufJeBedingung). Danach in beiden genau 5 min Durchsehen (Diff, Browser, Editor). (Der Plan-Satz "C startet nach 2 min Aufgabenblatt" ist verkürzt; der Vorher-Bogen liegt dazwischen.)
% - Je Person zwei gewertete Aufgaben, eine unter C, eine unter D: 2 Aufgaben x 2 Bedingungen x 2 Reihenfolgen (Plan, Zut; Zuteilung in 6.5).
% - Technik je Bedingung (Plan E5, Zut, Prüf, README des Werkzeugs; am 06.11. gegen die fünf Branches im Bare-Repo geprüft): Bare-Repo mit einem Branch je Aufgabe und Bedingung (studie-74-C, studie-74-D, studie-E-C, studie-E-D, studie-79-D), je ein Commit über dem Ausgangsstand mit .env, .env.local und .claude/settings.local.json. In allen: model claude-opus-5 (effort high aus den Benutzer-Settings), DISABLE_AUTOUPDATER=1, ENABLE_TOOL_SEARCH=false, CLAUDE_CODE_MCP_AUTO_BACKGROUND_MS=0, CLAUDE_CODE_MAX_MCP_DESCRIPTION_LENGTH=6000, MAX_MCP_OUTPUT_TOKENS=60000; eingefrorenes Binary 2.1.251 wie AP8. Nur in C: permissions.deny sperrt get_context_stepwise und get_context_interactive (geprüft 11.10. und 20.10.: /mcp zeigt 12 Werkzeuge, die Werkzeugliste des Modells hat 10; Prot 20.10.); dazu CODEBASE_GRAPH_NO_BROWSER=1, sonst öffnet get_context nach jeder Antwort den Neo4j Browser. Nur in D: CODEBASE_GRAPH_STEPWISE_UI=study (ohne Seed- und Rang-Panel, ohne Roh-Tab, Cap-Regler bleibt, Aufgabenzeile aufgeklappt; 4.5), EVAL_OUTPUT_DIR für den Export und REVIEW_TIMEOUT_MIN=15 (Leerlauffrist des Servers).
% - ENABLE_TOOL_SEARCH=false, weil Opus 5 am 09.10. get_context sonst ohne Schema aufrief (Aufgabentext wörtlich als task, kein userQuery; CLauf Runde 1, fünf von fünf Sitzungen laut README). Schon in 4.6 erklärt, hier nur Verweis.
% - Prompt je Aufgabe und Bedingung (eval/tasks/<Aufgabe>/prompts.md; es gibt keine eval/studie/prompts.md): C "Call the tool get_context from the codebase-graph server once with the task below, then solve the task in this repository." D "Call the tool get_context_stepwise ... once with the task below. Wait until I have confirmed the context in the browser, then solve the task ... using the confirmed context." Darunter der Aufgabentext. Start-Aufgabe im Terminal der Kopie legt ihn in die Zwischenablage, die Leitung fügt ein, der Proband drückt Enter, nicht Strg+Enter (Prot 21.10.). Wortlaut in den Anhang.
% - Wartezustand, Belege: README des Werkzeugs (Zeilen 322 bis 333), docs/tool-beschreibung.md, server.js, Commit 55731a1 (02.10.) für die 120-s-Beobachtung, Prot 11.10. Test 3 (Aufruf 13:39:15, Ergebnis 13:41:24, erste Dateilesung nach der Bestätigung). MCP_TOOL_TIMEOUT am 06.11. geprüft: in keinem Branch, nicht in den Benutzer-Settings, nicht in der Umgebung des Studienrechners gesetzt; empirische Untergrenze aus den Exporten: PG 74-D 22.10. 900 s, P03 E-D 03.11. 927 s, beide confirmed, timedOut false. REVIEW_TIMEOUT_MIN=15 in allen drei D-Branches; Leerlauf-Timer laut sitzung-anlegen.ps1 und 4.5.
%
% VERWEISE: 4.2 (Seeds ohne Eingriff), 4.5 (Werkzeuge, Studienseite, Export, Leerlauffrist), 4.6 (Werkzeugbeschreibung, ENABLE_TOOL_SEARCH), 6.4 (Vergleichswert, Transkript), 6.5 (Zuteilung, Ablauf).
% OFFEN: Version von Claude Code bei der 120-s-Beobachtung (vor dem 02.10., Prot hat keinen Eintrag). Erledigt 06.11.: Wartezeit der Oberfläche (15 min, Leerlauf), MCP_TOOL_TIMEOUT (nicht gesetzt, Voreinstellung, Sitzungen belegen über 15 min), Hintergrund-Pfad (Variable am 11.10. auf 2.1.251 geprüft, Befolgen der Anweisung ungeprüft, so geschrieben).

% WORTLAUT 6.2 (für den Text und Anhang app:aufgaben)
% - Prompt C, Aufgabe #74 (prompts.md, Wortlaut aus den Läufen 05.10. und 09.10.): "Call the tool get_context from the codebase-graph server once with the task below, then solve the task in this repository." Leerzeile, dann "Task: " und der Aufgabentext (6.3 WORTLAUT).
% - Prompt D, Aufgabe #74: "Call the tool get_context_stepwise from the codebase-graph server once with the task below. Wait until I have confirmed the context in the browser, then solve the task in this repository using the confirmed context." Leerzeile, "Task: " und der Aufgabentext. Bei E und #79 derselbe Rahmen mit dem jeweiligen Text; #79 nur D.
% - Start-Aufgabe (Funktion im Terminal der Kopie, kopie-terminal.ps1): legt den Prompt der Kopie in die Zwischenablage und startet das eingefrorene Binary; die Leitung fügt mit Strg+V ein, der Proband drückt Enter (Strg+Enter bricht den laufenden Schritt ab). Startbild muss zeigen: 2.1.251, Opus 5, keine IDE-Zeile; /mcp zeigt codebase-graph verbunden. Vertrauensfrage beim ersten Start mit Ja. In C zeigt /mcp trotzdem zwölf Werkzeuge, die Sperre wirkt in der Werkzeugliste des Modells (20.10. mit -p geprüft: zehn Werkzeuge, ohne Stepwise).
% - Ansage der Leitung beim Lauf in C: "Schau zu, was er tut. Denk laut." Kein Graph, der Agent arbeitet sofort.
% - Ansage beim Lauf in D: "Es öffnet sich der Browser mit dem Graphen. Du siehst, welche Stellen das Werkzeug dem Agenten geben will. Du kannst Haken abwählen, Nachbarn laden, Knoten hinzufügen, umsortieren. Wenn du fertig bist, drückst du 'Kontext bestätigen'. Dann arbeitet der Agent. Denk laut." Kuratierzeit von Enter bis Bestätigen; Grenze 15 min, bei 13 ansagen, bei 15 bitten zu bestätigen. Die "Restzeit" der Oberfläche ist ein Leerlauf-Timer, keine Gesamtgrenze.
% - Studienseite aus Sicht des Probanden (anleitung-proband.md, liegt im Sitzungsordner als 03-anleitung-oberflaeche.md und darf jederzeit geöffnet werden): links Suchtext ("Vektorisierter Suchtext"), "Knoten-Detail", Suche "Knoten hinzufügen"; Mitte Graph (Linien: Importe, Aufrufe, Vererbung, Funktion in Datei); rechts dieselben Knoten als Liste, darüber die Zahl der ausgewählten; oben "Restzeit", "Vorschau", "Zurück", "Weiter", darunter die Aufgabe zugeklappt. Stufe 1 Entitäten (x entfernt einen Begriff, Feld für eigene, "Deine Anfrage" mit "Übernehmen" zieht weitere Begriffe = retask); Stufe 2 Seed-Suche (Haken weg = Knoten fällt weg und lädt keine Nachbarn; "Verworfene zeigen", "+ Reaktivieren"); Stufe 3 Traversal ("+ Nachbarn laden" je Knoten, "+ Weiteren Hop für alle Knoten laden", "Letztes Laden zurücknehmen"); Stufe 4 Ranking ("Cap: Top-30" über der Liste, abgewählte Zeilen zählen nicht und zeigen "#–", Umsortieren durch Ziehen am Griff oder Klick auf die eigene Nummer "#…" (nicht "T#…"), Zielrang, Enter; "Kontext bestätigen"). Ab Stufe 2: Haken in der Liste; Klick auf Knoten zeigt links Quelltext, Verbindungen und zugeklappt die Begründung; die Karte erscheint schon beim Überfahren einer Listenzeile (daher views ab 300 ms, 6.4); "Knoten hinzufügen" nach Name oder Bedeutung, + am Treffer, auf Stufe 4 Abfrage des Rangs; "Vorschau" zeigt den Text, den der Agent bekäme; "Zurück" warnt, dass spätere Stufen neu gerechnet werden ("Verwerfen & zurück"; Abwahlen und von Hand hinzugefügte Knoten bleiben, Umsortieren und geladene Nachbarn nicht); "Restzeit" läuft bei Untätigkeit ab, jede Aktion setzt sie neu, danach bekommt der Agent die Auswahl, wie sie gerade ist. Die Zahlen "sim", "commits", "T#" und der Erklärtext unter dem Stufentitel sind für den Probanden ausdrücklich unwichtig. Schlusssatz der Anleitung: "Du entscheidest, ob und was du änderst. ... Sag dabei laut, was du denkst und warum du etwas tust."
% - Branches im Bare-Repo (studie\teamable-frontend.git), erwartete Commits Stand 14.10.: studie-74-C f7329ec, studie-74-D 449e244, studie-79-D 1778cba, studie-E-C f40e68e, studie-E-D 7090b5e (E-Branches seit 12.10. mit .env auf Backend-Port 8001, D-Branches seit 14.10. mit CODEBASE_GRAPH_STEPWISE_UI=study). Je Kopie ein Klon mit Arbeitsbranch <Pseudonym>-<Aufgabe>-<Bedingung>, node_modules aus vorlage-74 (für #74 und #79, gleiche yarn.lock) oder vorlage-E; kopien-pruefen.ps1 meldet verschobene Branches. Nach der Sitzung: git add -A, commit "<Pseudonym> <Aufgabe> <Bedingung>", push -u origin <Branch>; das Ergebnis bleibt als Branch.
% - Eingefrorenes Binary: %USERPROFILE%\claude-frozen\2.1.251\claude.exe (sha256-Anfang 8d1229a281281b98); claude aus dem PATH ist neuer und darf nicht laufen. VS-Code-Studienprofil mit eigenem Datenordner und Erweiterungsordner (nur drei Markdown-Erweiterungen, keine Claude-Code-Erweiterung, weil sich das Binary sonst als IDE verbindet und Werkzeuge bekommt, die es in AP8 nicht hatte).

Die Studie vergleicht zwei Bedingungen. In beiden löst derselbe Agent, Claude Code mit dem Modell Opus~5, dieselbe Aufgabe in einer eigenen Kopie der Codebasis und erhält seinen Kontext aus dem Graphen aus Kapitel~\ref{chap:werkzeug}. In C ruft er \texttt{get\_context} auf, bekommt den Kontext ohne Eingriff (Abschnitt~\ref{sec:abrufpipeline}), und der Proband sieht im Terminal zu. In D ruft er \texttt{get\_context\_stepwise} auf. Die Oberfläche aus Abschnitt~\ref{sec:betriebsmodi} öffnet sich in der Fassung für die Studie, im Folgenden die Studienseite. Der Proband kuratiert dort den Kontext Stufe für Stufe und bestätigt ihn, erst dann arbeitet der Agent. Alles andere ist gleich, auch der Prompt bis auf den Namen des Werkzeugs und den Satz, dass der Agent auf die Bestätigung warten soll (Anhang~\ref{app:aufgaben}). Warum der Vergleich nicht gegen Claude Code ohne Werkzeug läuft, begründet Abschnitt~\ref{sec:studienfragen}.

Sehen und Verändern lassen sich in D nicht trennen. Wer kuratiert, hat mehr Code gesehen, auch wenn er nichts verändert. Eine dritte Bedingung, in der der Proband nur liest, hätte bei acht Personen zwei bis drei je Zelle gelassen und einen Lesemodus sowie eine dritte Aufgabe verlangt. Ein Vorsprung von D gilt deshalb dem Verfahren als Ganzem, und die Kuratierzeit wird als Kosten berichtet (Abschnitt~\ref{sec:messung}). Aus dem stufenweisen Modus folgt auch der Vergleichswert für die erste Frage. Weil der Proband die Nachbarn selbst lädt, enthält ein Kontext ohne Eingriff nur die Seeds (Abschnitt~\ref{sec:konzeption}), und jedes Nachladen wäre gegen diesen Startzustand per Konstruktion eine Verbesserung. Ein fester Wert je Aufgabe taugt ebenfalls nicht, weil der automatische Kontext an der Anfrage des Agenten hängt (Abschnitt~\ref{sec:werkzeug-grenzen}). Vergleichswert ist deshalb je Person der Kontext, den \texttt{get\_context} mit den Eingabetexten aus dem Export derselben Sitzung geliefert hätte (Abschnitt~\ref{sec:messung}).

Jede gewertete Aufgabe läuft in beiden Bedingungen gleich ab (Abschnitt~\ref{sec:teilnehmende}): zwei Minuten Aufgabenblatt, Vorher-Bogen bis zur vierten Minute, dann schickt der Proband den vorbereiteten Prompt mit Enter ab. In C arbeitet der Agent sofort. In D kuratiert der Proband bis zur Bestätigung, höchstens 15 Minuten. Danach haben beide genau fünf Minuten für Diff, Browser und Editor. Jede Person bearbeitet eine Aufgabe unter C und eine unter D. Zwei Aufgaben, zwei Bedingungen und zwei Reihenfolgen ergeben vier Abfolgen zu je zwei Personen.

Die Einstellungen von Claude Code liegen in der Kopie der Codebasis: Ein lokales Bare-Repository hält je Aufgabe und Bedingung einen Branch mit einem Commit über dem Ausgangsstand, und jede Kopie wird daraus geklont. Überall gleich sind Modell, das eingefrorene Binary 2.1.251 (Abschnitt~\ref{sec:agentenverhalten}) und die Einstellungen aus Abschnitt~\ref{sec:werkzeug-grenzen}. Nur in C sind die beiden wartenden Werkzeuge gesperrt und der Neo4j Browser ist abgeschaltet, nur in D zeigt der Server die Studienseite und schreibt den Export.\footnote{In allen Branches: \texttt{claude-opus-5} mit \texttt{effort high}, \texttt{ENABLE\_TOOL\_SEARCH=false}, \texttt{CLAUDE\_CODE\_MAX\_MCP\_DESCRIPTION\_LENGTH=6000} (sonst kürzt Claude Code die Beschreibung auf 2048 Zeichen, und die Legende der JSON-Rückgabe fällt weg), \texttt{CLAUDE\_CODE\_MCP\_AUTO\_BACKGROUND\_MS=0}, automatische Updates aus. Nur in C: \texttt{permissions.deny} für \texttt{get\_context\_stepwise} und \texttt{get\_context\_interactive}, geprüft am 11.10.\ und 20.10.2026 mit 2.1.251 (die Werkzeugliste des Modells enthält zehn der zwölf Werkzeuge), und \texttt{CODEBASE\_GRAPH\_NO\_BROWSER=1}. Nur in D: \texttt{CODEBASE\_GRAPH\_STEPWISE\_UI=study}, \texttt{EVAL\_OUTPUT\_DIR} und \texttt{REVIEW\_TIMEOUT\_MIN=15}.} Ein Aufruf, der auf einen Menschen wartet, kann in Claude Code vor der Bestätigung enden. Nach 120 Sekunden verschiebt Claude Code ihn in den Hintergrund und lässt den Agenten weiterarbeiten. Die Studie schaltet das ab, und eine eigene D-Sitzung mit 2.1.251 zeigte, dass der Agent die erste Datei erst nach der Bestätigung las.\footnote{Das Verschieben wurde in der Vorbereitung vor dem 02.10.2026 beobachtet~[KLÄREN: Version von Claude Code]. Prüfung am 11.10.2026 mit 2.1.251 und Opus~5 auf Aufgabe \#74: Aufruf um 13:39:15, Ergebnis um 13:41:24, erste Dateilesung danach. Die Werkzeugbeschreibung bittet den Agenten zusätzlich, bei einem dennoch verschobenen Aufruf nicht zu arbeiten. Ob er das befolgt, ist ungeprüft, weil der Fall nicht mehr eintrat. Die Abbruchfrist \texttt{MCP\_TOOL\_TIMEOUT}, die die Fortschrittsmeldungen des Servers nicht verlängern, ist nicht gesetzt. Dass die Voreinstellung über 15 Minuten liegt, zeigen der Pilot vom 22.10.2026 (Aufruf 15:00 Minuten) und die Sitzung vom 03.11.2026 (15:27 Minuten), beide bestätigt.} Der Server selbst wartet in D 15 Minuten, als Leerlauffrist, die jede Aktion neu startet. Läuft sie ab, liefert er die letzte Auswahl als nicht bestätigt (Abschnitt~\ref{sec:betriebsmodi}). Die Kuratierzeit begrenzt also die Sitzungsleitung, nicht die Technik. Ob ein Agent vor der Bestätigung Dateien gelesen hat, zeigt im Nachhinein das Transkript von Claude Code (Abschnitt~\ref{sec:messung}).

\section{Aufgaben auf teamable}% 2 S.
% Inhalt: Referenzkontext, Spielraum je Aufgabe (13.10.).
\label{sec:aufgaben}
% ZIEL: Der Leser kennt die zwei gewerteten Aufgaben, die Auswahlregeln, den Referenzkontext und den Spielraum des Kuratierens je Aufgabe. Je Aufgabe der Text in einem Satz und eine Tabelle (Rolle, Ausgangsstand, Kern, Verständnisstellen, Deckung des automatischen Kontexts). Aufgabenblätter in den Anhang (app:aufgaben).
%
% ABWEICHUNG VOM EXPOSÉ, an dieser Stelle (Gl): Das Exposé nannte teamable-backend (Django, Python) und drei Aufgaben (Ex 03_methodik). Vorprüfung 05.09.: das Frontend trägt die Zwei-Kanten-Form, das Backend ist zu dünn (Tr19; ADR 0007, JSX-Komposition als CALLS, 4.3). Zwei gewertete Aufgaben plus Aufwärmen ist die vorab benannte Rückfallstufe (TR Regel 4, "Exposé 6.1").
%
% INHALT (geprüft)
% - Vier Aufgaben (Plan E4): #74 und E gewertet; #79 Aufwärmen, nur unter D (Logo und Kursname in der oberen Navigation, Stand 3777308); C (Suchfeld mit einem Klick leeren) Reserve, Text versiegelt, Stand 98ce1c8.
% - Auswahlregeln (TR, Stand 10.09., vorab festgelegt): Einstieg = Orte, die der Text in Prosa nennt (Route -> page.tsx, UI-Region -> Komponente). Gold-Kern gegen Nebenänderung. Verständnisstellen = zum Lösen nötig, nicht im Patch. Regel 4: (1) Text nennt kein Gold (strenger Schnitt aus eval/lib/named.js, UI-Strings gegen Kern-Hunks); (2) Kern mindestens zwei Quelldateien, getrennt ohne Einstiegsdateien; (3) mindestens ein Kern-Symbol oder eine Verständnisstelle, keine Nabe (unter 10 CALLS-Aufrufern), in Aufrufrichtung und ungerichtet mindestens zwei Kanten vom nächsten Einstieg. Strenge Lesart: nur dateiübergreifende Kanten zählen, Kern ohne Einstiegsdateien mindestens zwei Dateien. Beide gewerteten Aufgaben bestehen auch die strenge Lesart (T74, TE). Entspricht den Exposé-Eigenschaften 1 (unscharf, nennt Stellen nicht), 3 (zwei Kanten) und 4 (mehrere Dateien). Eigenschaft 2 (semantisch ähnliche, unbeteiligte Stellen in der Codebasis) wurde nicht geprüft; als Lücke nennen.
% - Ausgeschlossen (Plan, TR): #72 als Aufwärmaufgabe, weil es den Import-Hook aus #74 einbaut und den Zusammenhang vorwegnähme; #95 und #77 berühren dieselbe Seite; C teilt die Einstiegsseite mit E. Vor dem Scan: #73, #78 (Typänderungen), #85 (17 Dateien), #87 (Icon-Tausch), #100 (neue Seiten), #75, #76, #84, #93, #94 (eine Datei).
% - #74 (T74): echter PR, Ausgangsstand b768127, Gold-Commit 5ad8b7f (12.09.2024). Text nach Gutachten 10.09. entschärft (Lösungssatz zum React-Query-Hook, UI-String "Sections" und "without deleting its code" gestrichen; "Import students" bleibt, steht nur in der Einstiegsdatei). Einstieg StudentsPage. Kern drei Dateien (useStudents.tsx StudentsProvider, columns.tsx, page.tsx), ohne Einstiegsdatei zwei. Verständnisstellen vorab: useStudentsQuery (2 gerichtet, 1 dateiübergreifend), useStudents (2 gerichtet, 2 dateiübergreifend). Referenzkontext vom 12.09. (eigene blinde Bearbeitung, 40 min, erster Kontakt mit der Codebasis) ergänzt StudentTable, DataTableToolbar, DataTableViewOptions: zusammen 7 Referenzsymbole. Zulässige andere Orte: Sichtbarkeit in DataTable, invalidateQueries im Import-Hook. Defekt des PR: der Sections-Filter ruft getColumn weiter auf (console.error im Entwicklungsmodus), erst 70a5b93 behebt das; die Prüfliste führt das nur als Beobachtung.
% - E (TE): konstruiert, ohne PR, Ausgangsstand 98ce1c8 (HEAD), Text ist die dritte Fassung (konstruiert-entwuerfe-2026-10-05.md, zweite und dritte nach dem Retrieval-Ergebnis zur ersten). Einstiege TeamSetDetailPage, StudentsPageView. Kern aus der eigenen Lösung 3e17d65 (102 min): toaster.tsx Toaster, use-toast.ts reducer (TOAST_LIMIT 1 auf 10); reducer 4 gerichtet, 2 dateiübergreifend. Sieben Verständnisstellen (useExportTeamSetLms, useImportStudentsFromLms, useToast, toast, ToastProps, ToasterToast, addToRemoveQueue): zusammen 9 Referenzsymbole. Zulässige andere Orte: toast in use-toast.ts (Agentenlauf 05.10.), duration je Hook, toast.tsx (ToastClose, ToastViewport) als Ergänzung.
% - Referenzkontext je Aufgabe in reference-context.json: Herkunft (Datum, Methode, Dauer, Protokoll), Einstiege, Kern, Alternativen, Verständnisstellen, besuchte nicht nötige Stellen, externe Quellen, Suchbegriffe, priorExposure, verworfene Wege. Exposé: versioniert vor dem ersten Retrieval-Lauf, Betreuer beurteilt mindestens eine Aufgabe unabhängig. Stand: sealedIn in beiden task.md offen (seit 11.09. gilt statt Hash die Amendment-Regel: Änderungen mit Datum im Arbeitsprotokoll, TR); Zweitableitung durch den Betreuer nicht entschieden (TODO S3b). Plan stattdessen: die zwei Entwickler markieren nach ihrer Sitzung unabhängig die nötigen Stellen, bei der D-Aufgabe mit Vermerk. Bis 07.11. nicht geschehen (PG 22.10. ohne Markierung, P01 und P02 nicht gelaufen).
% - Spielraum 13.10. (eval/studie/kontext-gegen-referenz.js an den sechs C-Anfragen vom 09.10.; CLauf): Symbolebene #74 5, 6, 5 von 7 und E 8, 5, 7 von 9. Deckung (Symbolknoten oder Dateiknoten seiner Datei im Kontext) #74 dreimal 7 von 7, E 9, 8, 8 von 9; zweimal fehlt useExportTeamSetLms samt Datei. Alternativen #74 2 von 2, E 3, 2, 3 von 3. Folge: bei #74 Deckeneffekt, Kuratieren kann nur gleich bleiben oder verschlechtern; bei E kann Hinzufügen eine Datei ergänzen. Entwicklungswert aus zwei Stepwise-Exporten: eigene D-Sitzung zu E vom 05.10. Symbolebene 3 gegen 5 von 9, Deckung 9 gegen 9 bei 16 statt 30 Knoten (Plan E3, CLauf).
% - Grenzen (T74, TE): columns und TOAST_LIMIT sind Const-Deklarationen ohne Knoten; StudentsContextType, ToastProps, ToasterToast Typen ohne Kanten; toast hat im Graphen keinen Aufrufer. Die fünf Sekunden bei E sind die Voreinstellung von @radix-ui/react-toast 1.1.5 und stehen nicht im Repo (Kontrollitem 6.4). Liest man "red message in the corner" als UI-Region, wäre Toaster ein dritter Einstieg und der Kern ohne Einstiegsdateien eine Datei. Seeds hängen an der Anfrage des Agenten (E wörtlich: keine Kern-Datei in den Seeds; 4.6).
% - Umgebung je Aufgabe (Plan E5, Zut): #74 Backend-74 mit Canvas-Mock (Port 8765, je Import fünf neue Studierende, db.sqlite3.ausgang vor jeder Sitzung zurück), Frontend localhost:3000. E Backend-A ohne Mock, der Import muss scheitern, Frontend 127.0.0.1:3001 (Anmelde-Cookie token hängt am Hostnamen). Je Proband und Aufgabe ein eigener Klon aus dem Bare-Repo mit node_modules aus der Vorlage; das Ergebnis bleibt als Branch, kein Zurücksetzen.
%
% VERWEISE: 4.3 (Graph des Frontends, JSX als CALLS), 4.6 (Seeds an der Anfrage), 5.3 (Läufe 05.10. und 09.10. als Entwicklungsmessung), 5.6, 6.4 (Kontrollitems, Prüfliste), Anhang app:aufgaben.
% OFFEN: Markierung durch die Entwickler (07.11.: bisher nicht geschehen, im Text so gesagt; bleibt für P01 und P02 offen); Zweitableitung; sealedIn; Eigenschaft 2.
% Text am 06.11. geschrieben: sieben Absätze und Tabelle tab:aufgaben (beide Aufgaben als Spalten). Zweitableitung und sealedIn als Abweichung vom Plan geschrieben, Markierung durch die Entwickler als nicht geschehen (Marker am 07.11. aufgelöst), Eigenschaft 2 als Lücke. Zahlen aus reference-context.json, task.md und CLauf (Tabelle 13.10.).

% WORTLAUT 6.3 (Aufgabentexte für Anhang app:aufgaben, Referenzkontext, Akzeptanzkriterien)
% - Aufgabentext #74 (problem_statement, entschärft nach Gutachten 10.09.): "Two problems on the students page of a course: 1. The table shows a column with the sections a student belongs to. Sections are not populated for students yet, so the column is always empty and confuses users. Hide that column for now, in a way that makes it easy to bring back once sections are supported. 2. After using "Import students", the table keeps showing the old list until the browser page is reloaded. After a successful import the student list should refresh on its own." (Die zwei Punkte stehen als nummerierte Absätze.) Gestrichen gegenüber der Erstfassung: "The table data comes from a React Query hook behind a context provider; expose whatever is needed ...", der UI-String "Sections", "without deleting its code".
% - Aufgabentext E (dritte Fassung, 05.10.): "When an action fails, for example exporting a team set or importing students in a course that is not connected to Canvas, a red message appears in the corner and disappears by itself after about five seconds. If a second message comes up in the meantime, it replaces the first one instead of appearing next to it. Messages about a failure should stay until I close them, and when there are several messages they should be stacked so that all of them remain visible; other messages may still disappear on their own."
% - Aufgabentext #79 (Aufwärmen, PR #79 vom 23.09.2024, Ausgangsstand 3777308, Gold 4b5b29e): "The top navigation currently shows the plain text "Teamable Analytics" on the left. Replace it with the Teamable logo (an SVG file that will be placed under public/), and show next to it, separated by a vertical bar, the name of the course the user is currently in as an outlined badge. Clicking the logo goes to the home page. The navigation should be rendered only inside the layout of the course pages (the ones under a course id), not in the more general layout above it, because the course name is only known there." Einstieg Navbar; Kern Navbar (Logo in Link, Trennstrich, Badge mit courseName aus useCourse), Logo (SVG statt Text), course/layout.tsx (Navbar raus), course/[courseId]/layout.tsx (Navbar in CourseProvider); Verständnisstelle useCourse (Nabe, 14 Aufrufer). Keine Messaufgabe, weil Regel 4 nicht bestanden (Navbar zu Logo eine Kante, useCourse Nabe, UI-String "Teamable Analytics" steht direkt in Logo/index.tsx); fürs Aufwärmen erwünscht: kurze Wege, die Oberfläche zeigt schnell etwas Brauchbares. Vorbereitung: public/logo.svg als Platzhalter in der Kopie; kein Gate, keine Prüfliste; die App wird beim Aufwärmen nicht geöffnet (ob backend-74 zum Stand 3777308 passt, ist ungeprüft).
% - Aufgabentext C (Reserve, versiegelt, Stand 98ce1c8, eigene Bearbeitung in gate-C): "The search fields above the tables for team sets and projects have no quick way to reset them: users have to delete the text character by character. Add a one-click way to clear the search (and let the Escape key do the same), and the table must show all rows again right away." Teilt die Einstiegsseite mit E; bisher nicht gebraucht.
% - #74 Herkunft: Repo Teamable-Analytics/teamable-frontend, Ausgangs-Commit b768127 (PR #73 "render unassigned students", 12.09.2024), Gold-Commit 5ad8b7f (PR #74, 12.09.2024), 19 Commits vor HEAD 98ce1c8. Graph des Ausgangsstands: 96 Dateien, 172 Funktionen, 75 Klassen, 1038 Kanten. Einstieg "students page of a course" = src/app/(app)/course/[courseId]/students/page.tsx, StudentsPage. Kern (alles im Patch verlangt, keine Nebenänderung): students/(hooks)/useStudents.tsx StudentsProvider (reicht refetch der Query über den Context nach außen, StudentsContextType erweitert); students/(table)/columns.tsx (Const columns, kein Knoten; Spalte sections auskommentiert, Import Badge entfernt); students/page.tsx StudentsPage (Provider nach außen gezogen, neue StudentsPageView liest refetch aus useStudents und ruft es nach dem Import). Verständnisstellen: useStudentsQuery (refetch kommt von der React-Query-Abfrage; 2 gerichtet, 1 dateiübergreifend, 1 Aufrufer), useStudents (Context-Hook; 2 gerichtet, 2 dateiübergreifend über StudentsPage, StudentTable, useStudents; 4 Aufrufer). Referenzkontext (reference-context.json, Bearbeitung 12.09., 40 min, ohne Kenntnis des Patches) ergänzt StudentTable, DataTableToolbar, DataTableViewOptions. Zulässige andere Orte: DataTable in student-table.tsx (Sichtbarkeit), useImportStudentsFromLms (invalidateQueries im Hook). Auf HEAD ist die Lösung sichtbar (keine Sections-Spalte, Refresh nach Import); der PR-Autor wäre als Proband auszuschließen (nicht im Team). Hinweis im Kopf beider Bögen: Stand vom September 2024, weil der Import-Knopf auf HEAD seit dem PR innerhalb des Providers liegt (sonst wäre Frage 2 für Kenner von heute falsch gewertet).
% - #74 Akzeptanzkriterium (dreistufig): Gate tsc --noEmit und next lint grün (Voraussetzung 10.09. in gate-74 geprüft, Node 20.18.0, Yarn 1.22.22, Ausgangs- und Gold-Commit beide 0 Fehler, 0 Warnungen); Klickfolge 2a (keine Spalte "Sections", View-Menü bietet sie nicht an) und 2b (Import zeigt ohne Neuladen fünf Studierende mehr; Gegenprobe Ausgangsstand 11.10.: Tabelle bleibt stehen); sekundär Überschneidung mit dem Kern. Die Assertion aus task.md (a: keine Spalte "Sections" gerendert, b: nach erfolgreicher Import-Mutation zweiter Aufruf der Query-Funktion) wurde durch die Klickfolge der Prüfliste ersetzt. Beobachtung, kein Kriterium: TanStack 8.13.2 meldet "Column with id 'sections' does not exist" (Konsole, Entwicklungsmodus), sobald die Spalte fehlt und die Toolbar getColumn("sections") weiter aufruft; trifft den PR selbst, erst 70a5b93 behebt es.
% - E Herkunft: Ausgangs-Commit 98ce1c8 ("chore: add privacy policy page + cleanup for marketing pages (#100)"), kein Gold-Commit, Kern aus der eigenen Lösung 3e17d65 (loesung-autor.diff), Bearbeitung 05.10., 102 min, mit Protokoll. Graph: 103 Dateien, 185 Funktionen, 80 Klassen, 1160 Kanten (06.10.). Einstiege: "exporting a team set" = team-sets/[teamSetId]/page.tsx TeamSetDetailPage; "importing students" = students/page.tsx StudentsPageView. Kern: src/components/ui/toaster.tsx Toaster (variant aus den Props, duration ohne Ablauf bei destructive; kein gerichteter Weg, einziger Aufrufer RootLayout); src/hooks/use-toast.ts reducer, Fall ADD_TOAST, TOAST_LIMIT von 1 auf 10 (4 gerichtet, 2 dateiübergreifend über TeamSetDetailPage, useExportTeamSetLms, useToast, dispatch, reducer; 1 Aufrufer). Verständnisstellen: useExportTeamSetLms (onError löst die Meldung mit variant destructive aus), useImportStudentsFromLms (zweiter Auslöser, Hinweis und Fehlermeldung nacheinander), useToast (2 gerichtet, 2 dateiübergreifend, 9 Aufrufer, Nabenschwelle 10), toast (im Graphen ohne Aufrufer), ToastProps und ToasterToast (Typen ohne CALLS-Kante), addToRemoveQueue (TOAST_REMOVE_DELAY 8000 ms wirkt erst nach dem Schließen). Zulässige andere Orte: toast in use-toast.ts (duration nach variant plus Limit, eine Datei; Agentenlauf 05.10.); die fünf Hooks mit Fehlermeldung (use-export-team-set-lms, use-import-students-from-lms, use-import-student-gradebook-data, use-generate-teams, use-create-opt-in-quiz) je duration plus Limit; ToastClose und ToastViewport in toast.tsx als Ergänzung. Textfassungen: Fassung 1 vor jedem Retrieval, 2 und 3 nach dem Retrieval-Ergebnis zur ersten (konstruiert-entwuerfe-2026-10-05.md, Kandidat 1); Kandidaten 2 (F, Titel im Browser-Tab) und 3 (G, Kurs-Startseite nach dem letzten Schritt) erfüllten das Kriterium nicht; die Kandidaten wurden nach dem Retrieval ausgewählt, nicht davor (dort offen gesagt). Nennungsprüfung nach eval/lib/named.js für E nicht gelaufen (kein check.json, E steht nicht in make-tasks.js).
% - E Akzeptanzkriterium: Gate (06.10. geprüft, Node 20.18.0, Next 14.0.4, Ausgangsstand und eigene Lösung grün); Klickfolge: Manage Teams, Team-Set öffnen, Export oben rechts in einem Kurs ohne Canvas-Token (rote Meldung bleibt, mindestens 30 s warten); Seite Students, Import (Hinweis und Fehlermeldung gleichzeitig, der Hinweis verschwindet von selbst); sekundär Überschneidung mit Kern oder zulässigen Orten.
% - Auswahlregeln 1 bis 5 (eval/tasks/README.md, 10.09.): Regel 1 Einstieg = textbenannte Orte, nächster Einstieg zählt; bekannte Schwäche: Einstiege sind von der Nennungsprüfung ausgenommen, deshalb Kern ohne Einstiegsdateien getrennt. Regel 2 Gold-Kern gegen Nebenänderung (Umbenennung, Import-Sortierung, Formatierung, Stellen, die der Text nicht anspricht). Regel 3 Verständnisstellen vorab benannt, Naben zählen nicht. Regel 4 wie in INHALT; ehrliche Reihenfolge: Regel 4 (3) wurde nach dem ersten Scan von #72/#74/#88 von ungerichtet auf gerichtet gestellt (ungerichtet bestünden #72 und #79 über Rückwärtskanten und Naben); nach dem Gutachten 10.09. kamen Nabenschwelle, ungerichtete Mindestdistanz, strenge Lesart und UI-String-Abgleich dazu, danach alle sechs Instanzen neu gescannt. Regel 5 (vorregistriert): sealedIn-Commit, blinde Zweitableitung in zwei Stufen durch den Betreuer (nur Text, dann Diff; Skala 0 bis 2 für Unschärfe), erst dann Retrieval, Seed-Schicht hopSeed (Gold weder in den Seeds noch eine Kante davon entfernt, sonst "Gold in Seed-Reichweite", dann nur Kuratierungsfrage), PR-Autoren ausschließen, Vorfragebogen als Kovariate. Rollen vorab: #72 Aufwärmen (nach erstem Scan: Kern in der Einstiegsdatei), #74 und #88 Messaufgaben, #77/#95/#79 Kandidaten für die dritte. Stand jetzt: #88 nicht in der Studie (E an seiner Stelle), #79 Aufwärmen statt #72; sealedIn in beiden task.md offen; Zweitableitung nicht geschehen.
% - Rechnung 13.10. (kontext-gegen-referenz.js an den sechs C-Anfragen vom 09.10., drei je Aufgabe): Symbolebene #74 5, 6, 5 von 7; E 8, 5, 7 von 9. Deckung #74 7, 7, 7 von 7; E 9, 8, 8 von 9 (zweimal fehlt useExportTeamSetLms samt Datei). Alternativen #74 2 von 2; E 3, 2, 3 von 3. Früher (05.09., Default-Retrieval auf dem HEAD-Graphen zu einer Studierendenseiten-Aufgabe) lag der Import-Hook auf Rang 107; steht im Referenzkontext als priorExposure. Entwicklungswert E 05.10. (eigene D-Sitzung): Symbolebene 3 gegen 5 von 9, Deckung 9 gegen 9, 16 gegen 30 Knoten.
% - C-Läufe 09.10. mit Opus 5 (Entwicklungsmessung, Grundlage der Schlüssel): Arbeitszeit E 306 s, #74 153 s; Werkzeugaufrufe E 21 (1 get_context, 20 Bash), #74 17; Dateien gelesen E 15, #74 12; geändert E 4 plus Testdatei, #74 2; Gate danach beide grün; Eingabe an get_context nur task (Aufgabentext fast wörtlich), kein userQuery (Runde 1 ohne ENABLE_TOOL_SEARCH=false). Fable 05.10. (2.1.289): task kurz mit Symbolnamen, userQuery wörtlich. Dateien im Kontext (30 Knoten): E Opus 10, Fable 11; #74 Opus 14, Fable 18; Referenzdateien enthalten E 5 von 5, #74 6 von 6 in beiden.
% - Umgebung je Aufgabe (E5): #74 Backend-74 (Port 8000) mit Canvas-Mock (canvas-mock.py, Port 8765, je Import fünf neue Studierende), Frontend localhost:3000; E Backend-A (Port 8001) ohne Mock, Frontend 127.0.0.1:3001 (Anmelde-Cookie token hängt am Hostnamen); Neo4j 7687, Ollama 11434; Anmeldung in beiden Tabs mit dem Studienkonto; db.sqlite3.ausgang von Backend-74 wird nach jeder Sitzung zurückgesetzt (sitzung-nach.ps1). Je Proband drei Kopien (#79-D und die zwei gewerteten).
% GERÜST 6.3 (geschrieben 06.11., sieben Absätze plus tab:aufgaben, steht). Am 07.11. im Text: KLÄREN-Marker zur Markierung durch die Entwickler aufgelöst (nicht geschehen).

Die Aufgaben stammen aus teamable, einer Webanwendung zur Teambildung in Lehrveranstaltungen, an der die Teilnehmenden der Studie in ihrem Teamprojekt arbeiten (Abschnitt~\ref{sec:teilnehmende}). Bearbeitet wird das Frontend, in TypeScript mit React und Next.js geschrieben. Der ursprüngliche Plan sah das Backend in Python vor. Eine Vorprüfung ergab, dass dessen Graph zu dünn ist, um Aufgaben mit der unten beschriebenen Form zu tragen, während das Frontend sie über die Komposition von Komponenten hergibt, die der Graph als Aufrufkanten abbildet (Abschnitt~\ref{sec:codegraph}). Vier Aufgaben sind vorbereitet. Zwei werden gewertet: \#74, ein echter Pull Request aus der Geschichte des Projekts, und E, eine konstruierte Aufgabe auf dem aktuellen Stand. Dazu kommt \#79 zum Aufwärmen, immer unter D, mit Logo und Kursname in der oberen Navigation, sowie eine Reserveaufgabe mit versiegeltem Text, die bisher nicht gebraucht wurde. Der Plan nannte drei gewertete Aufgaben, zwei plus Aufwärmen war die vorab benannte Rückfallstufe. Die Aufgabenblätter stehen in Anhang~\ref{app:aufgaben}.

Die Regeln für die Auswahl wurden vor dem Durchsehen der Kandidaten festgelegt.\footnote{Stand 10.09.2026, im Werkzeug-Repositorium unter \texttt{eval/tasks/README.md}. Vor dem Scan ausgeschlossen wurden Pull Requests mit reinen Typänderungen, mit nur einer geänderten Datei, mit 17 Dateien, mit einem Icon-Tausch und mit neuen Seiten, die im Ausgangsgraphen fehlen. Nach dem Scan schieden \#72 als Aufwärmaufgabe aus, weil es den Import-Hook einbaut, um den es in \#74 geht, sowie \#95 und \#77, weil sie dieselbe Seite berühren.} Der Aufgabentext nennt nur das Verhalten, keine Datei, kein Symbol und keinen Bibliotheksnamen. Einstieg sind die Orte, die der Text in Prosa benennt, etwa eine Route oder eine Region der Oberfläche. Der Kern ist der Teil der Lösung, den der Text verlangt, ohne Nebenänderungen. Verständnisstellen sind Stellen, die man zum Lösen verstehen muss, die aber nicht im Patch stehen. Eine Aufgabe gilt als tragfähig, wenn der Kern ohne die Einstiegsdateien mindestens zwei Dateien umfasst und mindestens ein Kern-Symbol oder eine Verständnisstelle in Aufrufrichtung zwei Kanten über eine Dateigrenze vom nächsten Einstieg entfernt liegt, ohne eine Nabe mit zehn oder mehr Aufrufern zu sein. Beide gewerteten Aufgaben erfüllen das. Von den vier Eigenschaften, die der ursprüngliche Plan für Aufgaben verlangte, sind damit drei geprüft: Der Text nennt die Stellen nicht, die Lösung liegt zwei Kanten entfernt und verteilt sich über mehrere Dateien. Die vierte, semantisch ähnliche, aber unbeteiligte Stellen in der Codebasis, wurde nicht geprüft.

Aufgabe \#74 verlangt zweierlei auf der Studierendenseite eines Kurses: Eine stets leere Spalte soll so ausgeblendet werden, dass sie sich später leicht zurückholen lässt, und nach dem Import von Studierenden soll sich die Liste von selbst erneuern statt erst nach dem Neuladen der Seite. Der Text wurde nach einem Gutachten entschärft, ein Satz zum Lösungsweg und zwei Oberflächentexte, die direkt auf Kerndateien führten, sind gestrichen. Der Kern stammt aus dem Patch des Pull Requests: Der Provider der Seite reicht die Funktion zum erneuten Laden nach außen, die Spaltendefinition kommentiert die Spalte aus, und die Seite ruft das Laden nach dem Import auf. Aufgabe E verlangt, dass rote Fehlermeldungen, die nach etwa fünf Sekunden von selbst verschwinden und einander ersetzen, bis zum Schließen stehen bleiben und sich stapeln, während andere Meldungen weiter verschwinden dürfen. Sie hat keinen Patch. Der Kern ist die eigene Lösung des Autors in zwei Dateien, der Anzeigekomponente der Meldungen und dem Hook, der ihre Liste hält. Zwei weitere Lösungsorte, darunter einer in einer einzigen Datei, gelten vorab als zulässig. Tabelle~\ref{tab:aufgaben} stellt beide Aufgaben gegenüber.

\begin{table}[htbp]
  \centering
  \caption{Die zwei gewerteten Aufgaben. Deckung: Anteil der Referenzsymbole, deren Symbolknoten oder Dateiknoten im automatischen Kontext steht, gerechnet an drei Anfragen des Agenten je Aufgabe.}
  \label{tab:aufgaben}
  \small
  \begin{tabularx}{\textwidth}{l>{\raggedright\arraybackslash}X>{\raggedright\arraybackslash}X}
    \toprule
    & \#74 & E \\
    \midrule
    Herkunft & Pull Request vom September 2024, Text entschärft & konstruiert, ohne Patch, Text in dritter Fassung \\
    Ausgangsstand & 19 Commits vor dem aktuellen Stand & aktueller Stand \\
    Einstiege & Studierendenseite & Detailseite eines Team-Sets, Studierendenseite \\
    Kern & 3 Dateien, ohne Einstiegsdatei 2 & 2 Dateien, keine Einstiegsdatei \\
    Verständnisstellen & 5 & 7 \\
    Referenzsymbole & 7 & 9 \\
    Zulässige andere Orte & 2 & 3 \\
    Deckung je Anfrage & 7, 7, 7 von 7 & 9, 8, 8 von 9 \\
    Spielraum des Kuratierens & nur gleich oder schlechter & eine Datei nachholbar \\
    Umgebung & Backend mit lokalem Canvas-Mock, Import gelingt & Backend ohne Mock, Import scheitert \\
    \bottomrule
  \end{tabularx}
\end{table}

Der Referenzkontext je Aufgabe ist die Menge der Kern-Symbole und Verständnisstellen. Er entstand aus einer eigenen Bearbeitung des Autors auf dem Ausgangsstand, ohne Werkzeug und ohne Agenten, bei \#74 ohne Kenntnis des Patches und beim ersten Kontakt mit der Codebasis.\footnote{\#74 am 12.09.2026 in 40 Minuten, E am 05.10.2026 in 102 Minuten, jeweils mit Protokoll. Die Datei je Aufgabe hält neben Einstiegen, Kern, Alternativen und Verständnisstellen auch besuchte, aber nicht nötige Stellen, externe Quellen, Suchbegriffe, verworfene Wege und die Vorkenntnis des Autors fest.} Bei \#74 ergänzte die Bearbeitung drei Stellen, die der Patch nicht berührt, etwa die Tabelle, über die die Seite an die Spaltendefinition kommt. Der ursprüngliche Plan sah vor, dass der Betreuer mindestens eine Aufgabe unabhängig beurteilt. Das ist nicht geschehen. Stattdessen sollen die zwei Entwickler, die als Vergleichsfälle teilnehmen, nach ihrer Sitzung unabhängig die nötigen Stellen markieren (Abschnitt~\ref{sec:teilnehmende}). Bei der Generalprobe am 22.10.2026 ist das nicht geschehen. Bisher stammt der Referenzkontext aus einer Hand. Die Aufgabenblätter sind nicht mit einem Commit versiegelt, wie der Plan es vorsah. Seit dem 11.09.2026 gilt stattdessen, dass jede Änderung an Text, Einstiegen, Kern oder Verständnisstellen mit Datum und Grund im Arbeitsprotokoll steht.

Vor der ersten Sitzung wurde je Aufgabe gerechnet, welchen Teil des Referenzkontexts der automatische Abruf bereits liefert, und zwar an den sechs Anfragen, die der Agent in den Läufen unter C formuliert hatte (Abschnitt~\ref{sec:agentenverhalten}). Gezählt wird auf zwei Ebenen. Auf der Symbolebene zählt ein Referenzsymbol, wenn sein eigener Knoten im Kontext steht. Als gedeckt zählt es schon, wenn der Dateiknoten seiner Datei im Kontext steht, weil ein Dateiknoten dem Agenten die ganze Datei liefert. Die Symbolebene allein misst sonst die Körnung der Auswahl: In einer eigenen stufenweisen Sitzung zu E sank sie gegenüber dem automatischen Kontext von fünf auf drei von neun, während die Deckung bei neun von neun blieb, bei 16 statt 30 Knoten. Das Ergebnis steht in Tabelle~\ref{tab:aufgaben}. Bei \#74 deckt der automatische Kontext in allen drei Anfragen alle sieben Referenzsymbole, auch den Import-Hook, der in einer früheren Messung auf Rang 107 lag. Kuratieren kann den Kontext dort nur gleich lassen oder verschlechtern. Bei E fehlt in zwei von drei Anfragen der Hook, der den Export auslöst, samt seiner Datei. Dort kann Hinzufügen eine Datei nachholen. Deshalb wird die erste Frage je Aufgabe getrennt berichtet (Abschnitt~\ref{sec:auswertung}).

Vier Grenzen der Aufgaben sind vorab bekannt. Die Konstante, die bei E die Zahl der gleichzeitigen Meldungen begrenzt, und die Spaltendefinition bei \#74 sind Deklarationen ohne eigenen Knoten, und drei Typen haben Knoten ohne Kanten. Der Graph sieht einen Teil des Kerns also nicht. Die fünf Sekunden bei E sind die Voreinstellung einer Bibliothek und stehen nicht im Repositorium, was ein Kontrollitem in Abschnitt~\ref{sec:messung} nutzt. Liest man die rote Meldung in der Ecke bei E als Region der Oberfläche, wäre ihre Anzeigekomponente ein dritter Einstieg und der Kern ohne Einstiegsdateien nur noch eine Datei. Und welche Seeds der Abruf findet, hängt an der Anfrage, die der Agent formuliert: Mit dem Aufgabentext von E wörtlich als Anfrage liegt keine Kerndatei in den Seeds (Abschnitt~\ref{sec:werkzeug-grenzen}).


\section{Messung}% 2,5 S.
% Inhalt: Vorfragebogen, Fragen vorher und nachher, Erklären, UI-Export, Prüfliste für das Ergebnis des Agenten, TLX, SUS, Interview.
\label{sec:messung}
% ZIEL: Je Messgröße was, wann, wie gewertet, warum so. Tabelle "Messgrößen je Frage" (Frage, Instrument, Zeitpunkt, Wertung, Anhang). Bögen, Schlüssel, Aussagenlisten, Prüflisten in den Anhang (app:frageboegen); im Text je ein Beispielitem.
%
% WARUM VORHER UND NACHHER (Mess, Prot 08.10., Plan E3)
% - Fassung 02.10. (vorbereitete Diffs, Fragen nur nach dem Lauf) war aus Aufgabentext, Diff und Zusammenfassung des Agenten lösbar, in C und D gleich (sechs Durchsichten, Probelauf mit simulierten Probanden an den vier Läufen vom 05.10.). Optionen A bis H abgewogen, gewählt B (Vorher-Nachher mit Struktur-Items und Erklären, ohne Diffs) mit Auswertung im Stil von G (rein beschreibend). Bekannter Haken: der Vortest wird zum Suchauftrag; Gegenmittel sind nicht vorgefragte Items nachher und Interviewfrage 3.
% - Idee in einem Satz: In C und D steht am Ende dieselbe Art Zusammenfassung des Agenten; zweimal dasselbe fragen hebt das Lesen auf, übrig bleibt der Unterschied zwischen "im Terminal zusehen" und "im Graphen kuratieren" (Mess).
%
% INHALT (geprüft)
% - Vorfragebogen (eval/studie/vorfragebogen.md, Schlüssel vorfragebogen-schluessel.md, etwa 7 min): A Programmieren (Jahre; Handcode in den letzten drei Monaten; TypeScript, React, Next.js, Git je 1 bis 5), B teamable (Rolle, PR-Zahl, Beteiligung vor dem Teamprojekt, berührte Stellen, Nutzung im Browser), C KI-Werkzeuge (welche, wie oft, Erfahrung mit Agenten, Prüfverhalten), D Vorwissen zur Studie (schon gehört, Oberfläche gesehen), E drei kleine Aufgaben zum Codelesen mit Sicherheit. In der Arbeit nur je Gruppe berichtet (Plan, Einwilligung).
% - Vorher-Bogen (T74, TE vorher-bogen.md; nach dem Aufgabenblatt, vor dem Lauf, vor jedem Code): Kopf mit dem Aufgabentext im Wortlaut (seit 21.10.) und bei #74 dem Hinweis "Stand vom September 2024". Seit 02.11. zuerst Teil B Zusammenhänge (#74 Fragen 1 bis 6, E Fragen 1 bis 5; je drei Sachoptionen gleicher Länge plus "weiß ich nicht", richtige Antwort auf wechselnden Positionen, Sicherheit 1 bis 5), dann Teil A Ortsliste (#74 8 Stellen, E 7; spielt eine Rolle, keine Rolle, weiß ich nicht; beschreibend). Ebene Struktur in Alltagssprache, keine Datei- und Funktionsnamen. Kontrollitems: E Frage 5 (Herkunft der fünf Sekunden, kein Knoten, nicht im Diff) und #74 Frage 6 (Studierende je Seite, Const ohne Knoten; Grenze: "Rows per page 10" ist beim Durchsehen im Browser sichtbar, in D zeigt die Dateikarte die Zeile; je Person über views prüfen).
% - Nachher-Bogen (nachher-bogen.md): Teil B dieselben Items plus zwei nicht vorgefragte (#74 Fragen 7 und 8, E Fragen 6 und 7; Kontrolle gegen den Vortest als Suchauftrag); Teil C Arbeit des Agenten (#74 2 Fragen, E 3; beschreibend, Schlüssel je Sitzung aus Diff und Transkript; Lesekontrolle); Teil D TLX sechs Skalen 0 bis 100 in Fünferschritten mit Ankertexten (02.11.) und zwei Aussagen 1 bis 5 ("wusste, woran der Agent arbeitet", "könnte es erklären"); Teil A Ortsliste am Ende.
% - IM TEXT AUSSPRECHEN (Gutachten 2, 06.11.): Wiedererkennen gegen Verstehen. Teil A (Ortsliste) misst, ob die Person weiß, welche Stellen eine Rolle spielen; nur beschreibend. Teil B (Zusammenhänge) fragt, wie Stellen zusammenhängen und warum, in Alltagssprache ohne Datei- und Funktionsnamen; nur Teil B trägt die Hauptkennzahl. Das Erklären misst dasselbe frei formuliert. Damit misst die Studie Verständnis der Zusammenhänge, nicht das Wiedererkennen geänderter Dateien.
% - Wertung je Item: richtig 1, falsch 0, weiß ich nicht 0, keine Ratekorrektur (seit 12.10., schluessel.md beider Aufgaben). Minus ein Drittel im Mess-Entwurf ist überholt: bei drei Optionen blieb Raten mit +1/9 je Item im Vorteil (Prüf). Leere Items gelten als nicht erreicht (0, Zahl je Person daneben). Jedes Item einzeln berichtet.
% - Schlüssel je Aufgabe (schluessel.md): je Item Knoten, ob die Zusammenfassungen (Fable 05.10., Opus Runde 2 und 3 vom 09.10.) und der Diff es liefern, Status. Befunde: bei E liefert die Opus-Zusammenfassung mehr als Fable (Frage 2 variant destructive mit Datei und Zeile, Frage 5 Radix-Voreinstellung, 8000 ms nach dem Schließen); das Kontrollitem E5 fällt in jeder Sitzung, in der der Agent so zusammenfasst. Bei #74 hängt "geliefert" am Lösungsweg: Hook-Weg liefert Fragen 3 und 5, Umbau-Weg Fragen 1 und 2; die Items decken beide Wege ab, damit je Sitzung genug ungelieferte bleiben.
% - Zwei Schichten "geliefert": je Sitzung aus Transkript und Diff markiert, welche Items die Zusammenfassung beantwortet, vor dem Öffnen der Bögen (Leit "Nach der Sitzung"; seit P03 Dateien 12-/13-geliefert-*.md im Sitzungsordner). Regel: Beziehung genannt = ja, nur Datei- oder Symbolname = teils, nichts = nein (T74). Berichtet C gegen D bei gelieferten und bei nicht gelieferten Items (Plan).
% - Durchsehen genau 5 min (Diff, Browser, Editor); am Ende laut übernehmen, nachbessern oder verwerfen mit Grund, in den Mitschrieb; beschreibend (Ablauf).
% - Erklären 2 min, Tonaufnahme: Fenster der Kopie geschlossen, Chrome minimiert (seit 02.11.; schließen löscht die Anmeldung). Anweisung im Wortlaut mit Thema je Aufgabe (erklaeren.md: "Stell dir vor, jemand aus deinem Team arbeitet morgen an ... weiter"). Aussagenliste #74 A1 bis A10, E A1 bis A8, je Aussage Knoten, geliefert in den Läufen, Antwortoption im Bogen, Status. Zwei Werte je Person: übernommen (geliefert und genannt), eigenständig (nicht geliefert und genannt); Aussagen mit Antwortoption zählen nicht als eigenständig (Plan). Grundwert am Aufwärmen seit Pilot PG 22.10. gestrichen (die Zeile "Grundwert" in beiden erklaeren.md ist veraltet). Wertung durch die Sitzungsleitung aus dem Transkript; "ohne Hinweise auf die Bedingung" ist nominal, die Leitung kennt die Bedingung, keine verblindete Wertung (Grenze; Plan, Prüf). Seit 29.10. Verschriftung lokal (faster-whisper large-v3), Zeitleiste aus Transkript, Export und Claude-Transkript, Vorschlag eines lokalen Modells (qwen3:8b) je Aussage; die Wertung bleibt bei der Leitung (Ausw). Als KI-Nutzung in 1.6 kennzeichnen.
% - Eingabeweg der Bögen: der Proband liest über den Fernzugriff mit und sagt Antwort und Sicherheit, die Leitung trägt ein (21.10.: Klicks über RustDesk treffen die Kästchen nicht; 22.10.: Streuzeichen in den PG-Bögen). Entscheidung damit getroffen (Ablauf 02.11., Kopfzeile der Bögen). Effekt in C und D gleich, hebt sich innerhalb der Person auf; absolute Kalibrierung unter den Grenzen (Plan E5).
% - UI-Export (Plan E3, 4.5): je Stufe abgewählt, hinzugefügt, umsortiert, Übernahmen der Anfrage (retask); Dauer; toolOrder gegen order; views (Anzeigedauer der Detailkarte ab 300 ms, auch per Hover, Schwelle in der Auswertung, Kontrolle über die Aufnahme); ui: study. Kontextqualität auf Symbolebene und als Deckung gegen reference-context.json; Vergleichswert je Person neu gerechnet mit task und userQuery aus dem Export (kontext-gegen-referenz.js --export), nicht der Startzustand und kein fester Wert je Aufgabe (der automatische Kontext hängt an der Anfrage: 05.10. 1 bis 4 von 4 Referenzdateien je Anfrage). Dazu umgekehrter Anteil und Größe des Kontexts.
% - Prüfliste je Aufgabe (pruefliste.md; seit P03 als 10-/11-pruefliste-*.md im Sitzungsordner mit Screenshots): Gate tsc --noEmit und next lint (sitzung-nach.ps1 schreibt <Kennung>-gate-*.txt); Klickfolgen (#74 2a keine Spalte "Sections", View-Menü bietet sie nicht an; 2b Import zeigt fünf Studierende mehr ohne Neuladen; E Export- Meldung bleibt mindestens 30 s, beim Import Hinweis und Fehler gleichzeitig, Hinweis verschwindet); Überschneidung mit Kern und zulässigen Orten; aus dem Transkript Aufrufe, Arbeitszeit, Dateien gelesen und geändert, Lösungsweg, tsc oder lint selbst ausgeführt.
% - SUS (eval/studie/sus.md) am Ende, der Proband nennt je Aussage die Zahl. Interview mit vier Kernfragen im gleichen Wortlaut (Ablauf 02.11.), dabei Export der D-Aufgabe und Mitschrieb offen: 1 Moment ohne Verständnis; 2 woher wusste der Agent, wo er ändern muss (beide Aufgaben); 3 woran entschieden, was drinbleibt, nach den Fragen gesucht, falls nichts geändert warum; 4 aufgehalten oder gestört, bei welcher Art Aufgabe benutzen und bei welcher nicht, mit fester Nachfrage (Agent hat sich verlaufen, zweiter Versuch mit Oberfläche). Fragen 1 und 5 der Fassung 02.10. entfallen, Vorher-Nachher und Erklären messen sie (Plan).
% - Durchgehend lautes Denken; OBS mit Bildschirm, Mikrofon und Call-Ton (bisher eine gemischte Spur, Sprecher nachträglich zugeordnet; ab der nächsten Sitzung zwei Spuren geplant, Ausw); Transkript von Claude Code (sitzung-nach.ps1 kopiert es); Mitschrieb der Leitung (Zeiten je Teil, Kuratierzeit, Laufzeit, Urteil, Störungen).
% - Nicht eingeführt: ein Knopf je Knoten, der eine Erklärung von einem lokalen Modell holt (21.10. vertagt, hätte die Bedingung D verändert und Struktur-Items beantworten können; Prot 21.10.). D bleibt ohne.
% - Einschränkung: der Proband kennt beim Kuratieren die Items; Kuratieren ist auch Suche nach Antworten. Berichtet als "Kuratieren mit bekannten Fragen", Interviewfrage 3 fragt danach (Plan).
%
% VERWEISE: 4.5 (Export, Studienseite), 6.2 (Vergleichswert), 6.6 (Wertung), 7.2 bis 7.5, 8.2, Anhang app:frageboegen.
% OFFEN: Beispielitem für den Text wählen (Vorschlag E Frage 1 oder 3, #74 Frage 1); TLX-Auswertung benennen (Rohwerte ohne Gewichtung, Raw TLX) und zitieren (hart1988tlx, brooke1996sus in der Bib prüfen); Vertrauensskala jian2000trust aus dem Exposé wurde nicht erhoben, als Abweichung nennen; Auslese-Skript für die Bögen: gibt es nicht (07.11. geprüft), Bögen werden von Hand ausgelesen.

% GERÜST 6.4 (noch zu schreiben, 2,5 S.)
% Absatz 1: Was gemessen wird, in einem Satz je Größe, dann Tabelle tab:messgroessen (Spalten: Frage, Instrument, Zeitpunkt, Wertung, Anhang). Zeilen: Frage 2 Hauptkennzahl = Teil B der Bögen vorher und nachher (Zuwachs, nicht gelieferte Items); Frage 2 Erklären (übernommen, eigenständig); Frage 2 Kalibrierung (Sicherheit 1 bis 5); Frage 1 Export der Studienseite gegen neu gerechneten automatischen Kontext (Symbolebene, Deckung); E3 views, Eingriffe, lautes Denken, Interviewfrage 3; Nebenfrage Ergebnis = Prüfliste (Gate, Klickfolgen, Überschneidung); Nebenfrage Aufwand = Kuratierzeit, Raw TLX, SUS, Interview; Beschreibend = Teil A Ortsliste, Teil C, Urteil beim Durchsehen, Vorfragebogen je Gruppe.
% Absatz 2: Warum vorher und nachher (Fassung 02.10. verworfen; sechs Durchsichten, Probelauf mit simulierten Probanden; Optionen A bis H, gewählt B mit Auswertung wie G). Idee in einem Satz (Mess). Bekannter Haken Suchauftrag, Gegenmittel.
% Absatz 3: Aufbau der Bögen (Kopf mit Aufgabentext, Teil B zuerst, dann Teil A; nachher zusätzlich zwei nicht vorgefragte Items, Teil C, Teil D) und die Unterscheidung Wiedererkennen (Teil A) gegen Verstehen (Teil B, Erklären); ein Beispielitem im Wortlaut (Vorschlag: E Frage 1, Wortlaut unten). Wertung je Item, keine Ratekorrektur, leere Items.
% Absatz 4: Zwei Schichten "geliefert" mit der Regel (Beziehung = ja, Name = teils, nichts = nein) und dem Befund aus den Schlüsseln (bei E liefert die Opus-Zusammenfassung Fragen 2 und 5; bei #74 hängt es am Lösungsweg). Kontrollitems E5 und #74-6 mit ihren Grenzen.
% Absatz 5: Erklären (Anweisung im Wortlaut, Aussagenliste, zwei Werte, nicht verblindet, Transkription lokal, Vorschlag eines lokalen Modells, Wertung bei der Leitung; KI-Kennzeichnung nach 1.6).
% Absatz 6: Export der Studienseite und Vergleichswert (retask, views ab 300 ms, kontext-gegen-referenz.js --export); Transkript von Claude Code; Prüfliste; TLX und SUS in je einem Satz mit Zitat; Interview mit den vier Fragen; Eingabeweg der Bögen über die Leitung und was das für die Kalibrierung heißt; nicht eingeführter Erklärknopf; "Kuratieren mit bekannten Fragen".
% WORTLAUT 6.4, Vorfragebogen (eval/studie/vorfragebogen.md, etwa 7 min, Antworten sagt der Proband, keine richtigen Antworten in A bis D): A1 Jahre Programmieren (unter 1; 1 bis unter 3; 3 bis unter 6; 6 und mehr). A2 Code von Hand in den letzten drei Monaten (täglich oder fast täglich; etwa wöchentlich; seltener; gar nicht). A3 je Zeile 1 bis 5 für TypeScript, React (Komponenten, Hooks), Next.js, Git (Branch, Diff, Commit); Skala 1 kenne ich nicht, 2 schon gesehen, 3 kann ich lesen, 4 kann ich ohne KI ändern, 5 kann ich ohne KI neu schreiben. B1 Rolle im Teamprojekt (Frontend; Backend; Organisation, Tests, Dokumentation; nicht im Teamprojekt; mehrere möglich). B2 Pull Requests im Frontend (keine; 1 bis 3; 4 bis 10; über 10). B3 schon vor dem Teamprojekt am Code beteiligt (ja; nein). B4 selbst geänderte oder genauer angesehene Stellen (Anmeldung; Kursübersicht; Einrichtung eines Kurses; Studierendenseite; Team-Sets; Quiz-Seiten; obere Navigation; Meldungen unten rechts; Mobilansicht; keine). B5 teamable als Nutzer im Browser (regelmäßig; gelegentlich; einmal gesehen; nie). C1 KI-Werkzeuge in den letzten drei Monaten (Claude Code; GitHub Copilot; Cursor oder Windsurf; ChatGPT, Claude oder Gemini im Chat; andere; keine). C2 wie oft insgesamt (täglich; wöchentlich; seltener; gar nicht). C3 Erfahrung mit Agenten (noch nie; ein paar Mal ausprobiert; für einzelne Aufgaben regelmäßig; der größte Teil meines Codes entsteht so). C4 Prüfverhalten (Diff Zeile für Zeile; Diff überfliegen und ausprobieren; nur ausprobieren; meist ohne Prüfung). D1 schon etwas über die Studie gehört (nein; grob; Einzelheiten). D2 Oberfläche schon gesehen (nein; beim Vorgespräch kurz; selbst benutzt). E1 bis E3 drei kleine JSX-Aufgaben mit Sicherheit (E1 Counter mit useState, start 3, zweimal geklickt: a 5, b 3, c 2, richtig a; E2 Greeting ohne name: richtig c "Hallo Gast"; E3 useEffect mit [city]: richtig b "beim ersten Anzeigen und nach jeder Änderung von city"). Schlüssel: Summe 0 bis 3 als Niveau Codelesen; daneben beschreibend A1 bis A3, B1 bis B3, C3. Nur je Gruppe berichtet.
% WORTLAUT 6.4, Vorher-Bogen #74 (vorher-bogen.md, Kopf: "Sag die Antwort, die Sitzungsleitung setzt das Kreuz. Sicherheit: 1 heißt geraten, 5 heißt sicher. Aus dem Kopf antworten. 'Weiß ich nicht' ist immer eine gute Antwort, wenn du es nicht weißt. Der Code in dieser Sitzung ist ein Stand vom September 2024, vor späteren Umbauten."; dann der Aufgabentext). Teil B, je Frage a, b, c und d "Weiß ich nicht", Sicherheit 1 bis 5; richtige Antwort in Klammern: F1 "Woher bekommt die Tabelle die Liste der Studierenden?" a Die Tabelle selbst holt die Liste direkt vom Server; b Ein Baustein um die Tabelle herum holt und verteilt sie; c Die Seite holt die Liste und gibt sie der Tabelle mit (b; StudentsProvider, StudentTable, useStudents). F2 "Im Ausgangsstand des Codes: Wo liegt die Aktion, die der Knopf 'Import students' auslöst?" a In der Kopfzeile der Seite, außerhalb des Bausteins, der die Liste holt; b In der Leiste über der Tabelle, neben der Suche; c Innerhalb des Bausteins, der die Liste holt und an die Tabelle gibt (a; StudentsPage ruft useImportStudentsFromLms außerhalb des Providers). F3 "Wird der Import der Studierenden noch an einer anderen Stelle der Anwendung benutzt?" a Ja, auf der Seite mit den Team-Sets; b Nein, nur auf der Studierendenseite; c Ja, auch im Einrichtungsassistenten des Kurses (c; useSetupSteps). F4 "Der Sections-Filter in der Leiste über der Tabelle: was passiert mit ihm, wenn die Spalte aus der Spaltenliste verschwindet?" a Er bleibt stehen und zeigt eine leere Auswahl; b Er verschwindet mit, weil er prüft, ob die Spalte da ist; c Er muss von Hand entfernt werden (b; DataTableToolbar, table.getColumn("sections") &&). F5 "Womit ist die Anfrage an den Server für die Studierendenliste gekennzeichnet?" a Mit dem Kurs und den Einstellungen der Tabelle (Seite, Sortierung, Suche); b Nur mit der Kennung des Kurses; c Mit der angemeldeten Person und dem Kurs (a; useStudentsQuery, queryKey mit queryString). F6 (Kontrolle) "Wie viele Studierende zeigt die Tabelle beim ersten Öffnen je Seite?" a 25; b 10; c 50 (b; INITIAL_PAGE_SIZE = 10, Const ohne Knoten, in keinem Diff und keiner Zusammenfassung; Grenze: "Rows per page 10" steht beim Durchsehen im Browser, in D zeigt die Dateikarte von useStudents.tsx die Zeile, je Person über views prüfen). Teil A "Welche dieser Stellen im Projekt spielen eine Rolle dabei, welche Spalten die Tabelle zeigt und wann sie ihre Daten neu lädt?" (spielt eine Rolle; keine Rolle; weiß ich nicht): 1 Die Liste, die festlegt, welche Spalten die Tabelle hat (Rolle); 2 Die Stelle, an der die Studierendenliste vom Server geholt wird (Rolle); 3 Die Leiste über der Tabelle mit Suche, Filter und View-Menü (Rolle); 4 Der Baustein, der den Import ausführt (Rolle); 5 Die Meldungen unten rechts (keine; beim Hook-Weg vertretbar, nur beschreibend); 6 Die obere Navigation (keine); 7 Die Seite mit den Team-Sets (keine, Distraktor mit Nähe); 8 Die Anmeldung (keine).
% WORTLAUT 6.4, Nachher-Bogen #74: Kopf wie vorher mit "nicht nachsehen", Aufgabentext, Teil B Fragen 1 bis 6 identisch, dazu F7 "Das View-Menü über der Tabelle: woher weiß es, welche Spalten es anbietet?" a Aus einer eigenen Liste im Menü; b Vom Server zusammen mit der Liste; c Aus der Spaltenliste der Tabelle, alle ausblendbaren Spalten (c; DataTableViewOptions, getCanHide). F8 "Die Einstellungen, die die Tabelle gerade hat: wo werden sie gehalten?" a In der Tabelle selbst, getrennt von den Daten; b Im selben Baustein, der auch die Liste holt; c Auf dem Server je Person (b; useStudentsFilters im Provider; a halb wahr, weil DataTable Sortierung und Sichtbarkeit lokal hält). Teil C: F9 "An welchen dieser Stellen hat der Agent etwas geändert? (mehrere möglich)" a Spaltenliste; b Baustein, der die Liste holt; c Seite, Kopfzeile mit Import-Knopf; d Baustein, der den Import ausführt; e Weiß ich nicht (Schlüssel je Sitzung aus dem Diff: Hook-Weg a, d; Umbau-Weg a, b, c). F10 "Hat der Agent sein Ergebnis geprüft, bevor er fertig war (Typprüfung, Lint)?" Ja; Nein; Weiß ich nicht (aus dem Transkript). Teil D "Wie war die Aufgabe für dich": je Zeile 0 bis 100 in Fünferschritten, hohe Zahl heißt überall schwerer: Geistige Anforderung (wie viel Denken, Entscheiden, Suchen war nötig; 0 wenig, 100 sehr viel); Körperliche Anforderung (0 keine, 100 sehr hoch); Zeitliche Anforderung (Zeitdruck; 0 keinen, 100 sehr starken); Leistung (0 sehr gut gelungen, 100 misslungen); Anstrengung (wie hart musstest du arbeiten; 0 gar nicht, 100 sehr hart); Frustration (wie unsicher, entmutigt, gereizt; 0 gar nicht, 100 sehr). Dann zwei Aussagen 1 (stimme nicht zu) bis 5 (stimme voll zu): "Ich wusste während des Laufs, woran der Agent gerade arbeitet"; "Ich könnte jemandem aus dem Team erklären, was geändert wurde und warum". Teil A am Ende wie vorher. Positionen der richtigen Antworten #74: b, a, c, b, a, b, c, b; Optionslängen je Frage innerhalb von zehn Zeichen.
% WORTLAUT 6.4, Vorher-Bogen E (Kopf ohne den Satz zum Stand 2024): F1 "Wird die Fehlermeldung beim Export auf der Seite selbst erzeugt oder woanders?" a In einem Baustein, den die Seite benutzt; b Auf der Seite selbst, direkt beim Export-Knopf; c Auf dem Server, der die Antwort zurückschickt (a; TeamSetDetailPage ruft useExportTeamSetLms; trennte im Probelauf D 4 von 4, C 1 von 4). F2 "Wo wird festgelegt, dass eine Meldung rot ist?" a Dort, wo alle Meldungen angezeigt werden; b Dort, wo die Meldung ausgelöst wird; c Dort, wo die Liste der Meldungen geführt wird (b; variant destructive im Hook; bei Opus 09.10. geliefert). F3 "Beim Import erscheinen ein Hinweis und danach eine Fehlermeldung. Woher kommen die beiden?" a Beide aus der Stelle, die alle Meldungen anzeigt; b Der Hinweis von der Seite, die Fehlermeldung vom Server; c Beide aus der Stelle, die den Import ausführt (c; zwei toast-Aufrufe in useImportStudentsFromLms). F4 "Woher kommt der Wortlaut der roten Meldung beim Import?" a Der Server schickt den Text mit; b Der Text steht fest an der Stelle, die den Import ausführt; c Der Text steht an der Stelle, die alle Meldungen anzeigt (b; fester Text im onError). F5 (Kontrolle) "Wo ist festgelegt, dass eine Meldung nach etwa fünf Sekunden verschwindet?" a An einer Stelle im Projekt, als feste Zahl für alle Meldungen; b Nirgends im Projekt, der eingebundene Baustein bringt die Dauer mit; c An jeder Stelle, die eine Meldung auslöst, einzeln (b; Voreinstellung von @radix-ui/react-toast 1.1.5, kein Knoten; Opus 09.10. nennt die Herkunft in der Zusammenfassung, dann fällt die Kontrolle). Teil A "Welche dieser Stellen spielen eine Rolle dabei, wie Fehlermeldungen erscheinen und verschwinden?": 1 Die Stelle, die eine neue Meldung anlegt und die Liste der offenen Meldungen führt (Rolle); 2 Die Festlegung, wie eine Meldung aussieht (Farbe, Schließen-Knopf) (Rolle); 3 Die Stelle, die beim Anmelden die Fehlermeldung auslöst (Rolle; zeigt den Wortlaut-Effekt der Zusammenfassung); 4 Die obere Navigation (keine); 5 Die Festlegung der Spalten der Studierendentabelle (keine); 6 Der Server (Backend) (keine; einziger Distraktor mit Streuung); 7 Die Suche in der Studierendentabelle (keine).
% WORTLAUT 6.4, Nachher-Bogen E: Teil B 1 bis 5 identisch, dazu F6 "Wo ist die Anzeige der Meldungen in die Anwendung eingehängt?" a In jeder Seite einzeln; b In der oberen Navigation; c Im äußeren Rahmen, der alle Seiten umschließt (c; RootLayout, Toaster). F7 "Beim Import erscheint zuerst ein Hinweis. Wann wird er ausgelöst?" a Bevor die Anfrage an den Server abgeschickt wird; b Sobald der Server die Anfrage angenommen hat; c Erst wenn die Antwort da ist (a; toast vor mutateAsync). Teil C: F8 "An welchen dieser Stellen hat der Agent etwas geändert? (mehrere möglich)" a Dort, wo alle Meldungen angezeigt werden; b Dort, wo eine Meldung erzeugt und in die Liste aufgenommen wird; c In der Festlegung von Aussehen und Schließen-Knopf; d Dort, wo Export oder Import die Meldung auslösen; e Weiß ich nicht (aus dem Diff; 05.10. C-Lauf a, b, c; D-Lauf b, c; eigene Lösung a, b). F9 "Ist der Knopf zum Schließen bei roten Meldungen jetzt zu sehen, ohne dass man mit der Maus darüberfährt?" Ja; Nein; Weiß ich nicht, mit Sicherheit (aus dem Diff; 05.10. in allen vier Läufen Ja). F10 Agent hat geprüft, wie #74. Teil D und Teil A wie #74. Positionen der richtigen Antworten E: a, b, c, b, b, c, a.
% WORTLAUT 6.4, Beispielitem für den Text (Vorschlag E Frage 1, weil es im Probelauf am stärksten trennte und die Ebene "Struktur in Alltagssprache" zeigt): "Wird die Fehlermeldung beim Export auf der Seite selbst erzeugt oder woanders? (a) In einem Baustein, den die Seite benutzt. (b) Auf der Seite selbst, direkt beim Export-Knopf. (c) Auf dem Server, der die Antwort zurückschickt. (d) Weiß ich nicht. Sicherheit 1 bis 5." Richtig a; im Graphen die Kante von der Seite zum Hook. Alternative #74 Frage 1 (Provider-Muster).
% WORTLAUT 6.4, Schlüssel (schluessel.md #74 Stand 11.10., E Stand 09.10.): Wertung richtig 1, falsch 0, weiß ich nicht 0, keine Ratekorrektur (12.10.); leer = nicht erreicht (0 in der Hauptkennzahl, Zahl je Person daneben); jedes Item einzeln; Hauptkennzahl = Zuwachs auf den Fragen 1 bis 5, soweit in der Sitzung nicht geliefert; die nur nachher gefragten Items (#74 F7, F8; E F6, F7) getrennt. Geliefert-Spalten aus den Läufen: #74 Fable 05.10. und Opus Runde 2 (Hook-Weg: Fragen 3 und 5 geliefert, 4 immer), Opus Runde 3 (Umbau-Weg: Fragen 1, 2, 4, teils 7 und 8); E Fable C und D 05.10. (nur Frage 2 teils), Opus Runde 2 (Fragen 2 und 5 geliefert, 1 indirekt; 3, 4, 6, 7 ungeliefert). Status-Spalten im Schlüssel (bleibt, prüfen, kann raus) stammen aus der Planung vor dem Pilot; gestrichen wurde nichts, alle Items blieben in den Bögen. Teil D: die zwei Aussagen "können raus" (Probelauf C und D gleich 2 bis 3; Selbsteinschätzung), blieben aber drin. Offen im Schlüssel: Prüfung der Fragen durch jemanden, der nur Aufgabentext und Zusammenfassung kennt (nicht geschehen); Kurzpilot (entfallen).
% WORTLAUT 6.4, Erklären (erklaeren.md): Anweisung #74 "Stell dir vor, jemand aus deinem Team arbeitet morgen an der Studierendentabelle weiter. Erklär ihm in zwei Minuten: Was wurde geändert, warum war es vorher so, und worauf muss er achten. Es gibt kein richtig oder falsch, sag einfach, was du weißt." Bei E "an den Meldungen weiter", sonst gleich. Ablauf: Fenster der Kopie geschlossen, Chrome minimiert, Tonaufnahme läuft, zwei Minuten. Aussagenliste #74 (Knoten; Antwortoption im Bogen): A1 Die Spalten stehen in einer eigenen Liste, die Tabelle bindet sie ein (columns.tsx zu StudentTable; keine; in jeder Zusammenfassung teils). A2 Der Sections-Filter prüft, ob die Spalte existiert, und verschwindet mit ihr (DataTableToolbar; F4; immer geliefert). A3 Das View-Menü bietet jede ausblendbare Spalte an; nur ausblenden reicht nicht, wenn das Menü sie zurückholt (DataTableViewOptions; F7). A4 Die Liste wird von einem umschließenden Baustein geholt und weitergegeben, die Tabelle holt sie dort ab (StudentsProvider, useStudents; F1; nur Umbau-Weg geliefert). A5 Der Import-Knopf liegt außerhalb dieses Bausteins und kommt deshalb nicht an die Nachladefunktion (StudentsPage; F2; nur Umbau-Weg). A6 Der Schlüssel der Liste enthält die Filter; Neuladen muss alle Varianten treffen (useStudentsQuery; F5; nur Hook-Weg). A7 Der Import wird auch vom Einrichtungsassistenten benutzt, eine Änderung im Hook wirkt dort mit (useSetupSteps, useImportStudentsFromLms; F3; nur Hook-Weg). A8 Bei Fehlschlag zeigt der Import nur eine Meldung, neu geladen wird nur bei Erfolg (onError; keine Option). A9 Die Abfrage der Liste startet erst, wenn die Filter aus der Adresse gelesen sind (enabled: !initialFiltersLoading; keine; in keiner Zusammenfassung). A10 Suche, Seite und Sortierung werden in die Adresse gespiegelt (useStudentsFilters, replaceState; keine). "Eigenständig" läuft bei #74 über A8 bis A10 plus A3 bis A7 je nach Lösungsweg. Aussagenliste E: A1 Fehlermeldungen werden in den Bausteinen der Aktionen (Hooks) ausgelöst, nicht auf der Seite (useExportTeamSetLms, useImportStudentsFromLms; F1). A2 Alle Meldungen laufen über eine gemeinsame Stelle mit der Liste und werden an einer Stelle angezeigt (useToast, reducer, Toaster; F6 teils). A3 Die Liste war auf eine Meldung begrenzt, deshalb ersetzte die zweite die erste (reducer, TOAST_LIMIT; keine; immer geliefert). A4 Die fünf Sekunden sind die Voreinstellung des Bausteins, nicht im Projekt festgelegt (kein Knoten; F5). A5 Rot ist eine Art von Meldung, daran hängt das Stehenbleiben (variant destructive; keine; immer geliefert). A6 Die Wartezeit im Projekt (8000 ms) wirkt erst nach dem Schließen (addToRemoveQueue; keine; bei Opus geliefert). A7 Eine zentrale Änderung gilt für alle Fehlermeldungen, auch Anmelden; eine Änderung nur am Export erfasst den Import nicht (Aufrufer von useToast; keine; immer geliefert). A8 Beim Import kommen Hinweis und Fehlermeldung nacheinander aus derselben Stelle (useImportStudentsFromLms; F3 und F7). Wertung: Transkript ohne Hinweise auf die Bedingung lesen (Wörter wie Graph, Knoten, ausgewählt, Karte entfernen), je Aussage genannt oder nicht; Nennung heißt, die Beziehung ist erkennbar, auch in anderen Worten. Beispiel A4 (#74) genannt: "die Tabelle kriegt die Liste von dem Ding drumherum, das sie lädt"; nicht genannt: "useStudents wurde geändert". Beispiel A1 (E) genannt: "die Meldung kommt aus dem Hook für den Export, nicht von der Seite". Die Zeile "Grundwert" in beiden erklaeren.md und die Datei zu #79 (G1 bis G4, eine Minute, "an der oberen Navigation weiter") sind seit 23.10. außer Kraft.
% WORTLAUT 6.4, Prüflisten (pruefliste.md, füllt die Leitung nach der Sitzung aus, Screenshots seit P03): Kopf Pseudonym, Bedingung, Branch. 1 Gate: tsc --noEmit ohne Fehler; next lint ohne Fehler und Warnungen (sitzung-nach.ps1 schreibt <Kennung>-gate-<Aufgabe>-<Bedingung>.txt). 2 Klickfolge #74 (Entwicklungsmodus, Konsole offen): 2a Students-Seite zeigt keine Spalte "Sections", View-Menü bietet sie nicht an; Beobachtung Konsolenfehler (ja/nein, Wortlaut) und ob die Toolbar berührt wurde; 2b "Import students" mit laufendem Mock: Meldung "Imported students", Seite bleibt bedienbar, Tabelle zeigt ohne Neuladen fünf mehr; Beobachtung, ob das Nachladen an onSuccess oder an await hängt und nur bei Erfolg greift. Klickfolge E: Export in Kurs ohne Canvas-Token, rote Meldung erscheint unten rechts und bleibt mindestens 30 s; Import: Hinweis und Fehlermeldung gleichzeitig, der Hinweis verschwindet. 3 Tabelle geänderte Dateien gegen Kern, zulässiger anderer Ort, keins von beiden. 4 Aus dem Transkript: Aufrufe von get_context bzw. get_context_stepwise, Arbeitszeit vom Prompt bis zur letzten Antwort (s), Dateien gelesen, Dateien geändert, Lösungsweg (#74: Hook, Umbau der Seite, Sichtbarkeit in der Tabelle, anderes), Typprüfung oder Lint selbst ausgeführt. Ergebnis: bestanden (Gate und Klickfolgen vollständig) oder nicht bestanden mit Grund.
% WORTLAUT 6.4, SUS (eval/studie/sus.md): Kopf "Gemeint ist die Oberfläche im Browser, in der du die Auswahl der Codestellen für den Agenten gesehen und verändert hast: die vier Stufen mit Graph und Liste, bei der Aufwärmaufgabe und bei der Aufgabe, bei der sich der Browser geöffnet hat. Nicht gemeint sind Claude Code im Terminal und die teamable-Anwendung." Je Zeile 1 (stimme gar nicht zu) bis 5 (stimme voll zu), aus dem Bauch, bei Unsicherheit 3. Die zehn Aussagen in der Standardreihenfolge nach Brooke (1 regelmäßig benutzen; 2 unnötig komplex; 3 einfach zu benutzen; 4 Hilfe einer technisch erfahrenen Person nötig; 5 Funktionen gut aufeinander abgestimmt; 6 an zu vielen Stellen uneinheitlich; 7 die meisten lernen sehr schnell; 8 sehr umständlich; 9 sehr sicher im Umgang; 10 viele Dinge lernen müssen), Wortlaut "Ich denke, dass ich diese Oberfläche gern regelmäßig benutzen würde" usw. Auswertung nach Brooke: ungerade Items Wert minus 1, gerade 5 minus Wert, Summe mal 2,5 (0 bis 100). Bib prüfen: brooke1996sus.
% WORTLAUT 6.4, TLX: Rohwerte je Skala 0 bis 100, ohne Gewichtungspaare, also Raw TLX; Ankertexte seit 02.11. in den Bögen (oben). Bib prüfen: hart1988tlx; für Raw TLX ggf. Hart 2006 ("NASA-Task Load Index (NASA-TLX); 20 years later") ergänzen. Vertrauensskala jian2000trust aus dem Exposé wurde nicht erhoben; als Abweichung in einem Satz.
% WORTLAUT 6.4, Interview (ablauf-vorlage.md, 02.11.; Export der D-Aufgabe und Mitschrieb offen; Stichworte je Frage in den Mitschrieb): 1 "Gab es einen Moment, in dem du nicht wusstest, was der Agent gerade tut oder warum? Was hast du dann gemacht?" 2 "Woher wusste der Agent bei der ersten Aufgabe, wo er ändern muss? Und bei der zweiten?" 3 "Du hast in der Auswahl [aus dem Export] geändert. Woran hast du entschieden, was drinbleibt? Hast du nach den Fragen vom Bogen gesucht?" Falls nichts geändert wurde: "Warum?" 4 "Was hat dich an der Oberfläche aufgehalten oder gestört? Bei welcher Art Aufgabe würdest du sie benutzen, bei welcher nicht?" Nachfrage, falls nicht von selbst genannt: "Stell dir vor, der Agent hat eine Aufgabe schon einmal falsch gelöst oder sich in der Codebasis verlaufen. Würdest du es beim zweiten Versuch mit der Oberfläche machen? Warum?" Zum Schluss: "Bitte bis zum letzten Termin nichts über Aufgaben und Fragen weitergeben."
% WORTLAUT 6.4, Durchsehen (Ansage): "Der Agent ist fertig. Du hast fünf Minuten: Schau dir an, was er geändert hat. Du kannst den Diff lesen, die App im Browser ausprobieren, im Editor nachsehen. Am Ende sagst du mir: Würdest du das übernehmen, nachbessern oder verwerfen, und warum." Hinweis der Leitung je Aufgabe: #74 Import im Browser zeigt, ob die Tabelle ohne Neuladen um fünf wächst; E Fehlermeldungen auslösen (Import ohne Canvas zweimal kurz hintereinander oder Team-Set exportieren), bleiben sie stehen und stapeln sie sich.
% WORTLAUT 6.4, Export der Studienseite (Plan E3, 4.5): JSON je Sitzung unter studie\export, nach der Sitzung nach <Pseudonym>-export; Felder: je Stufe abgewählt, hinzugefügt, umsortiert, Übernahmen der Anfrage (retask), Dauer, toolOrder gegen order, views (Anzeigedauer der Detailkarte, auch per Hover), ui: study, confirmed, timedOut, task und userQuery des Agenten. Vergleichswert: kontext-gegen-referenz.js --export rechnet get_context mit task und userQuery aus dem Export neu und vergleicht beide gegen reference-context.json (Symbolebene, Deckung, umgekehrter Anteil, Größe).
% WORTLAUT 6.4, Auswertung der Aufnahme (studie\auswertung\README.md, 29.10.): OBS-Video, ffmpeg, faster-whisper large-v3 lokal (etwa 15 min je 100 min Aufnahme), Sprecherzuordnung nachträglich über Stimmabdrücke (sprecher-zuordnen.py), ab der nächsten Sitzung zwei Tonspuren (Mikrofon = Leitung, Discord = Proband). Ergebnis je Person: transkript-lesbar.md (L und P, Uhrzeit, Abschnitt), transkript.txt, zeitleiste.md (Transkript verzahnt mit Enter, Aktionen der Oberfläche, angesehenen Knoten, Werkzeugaufrufen und Sätzen des Agenten), je Abschnitt eine Textdatei, je Aufgabe erklaeren-vorschlag.md (qwen3:8b schlägt vor, welche Aussagen genannt wurden; Wertung bleibt bei der Leitung; --ohne-ollama schaltet ab), interview-vorschlag.md (Zusammenfassung je Frage mit Zitaten). Abschnittsgrenzen aus Enter, Ende des Laufs und Werkzeugaufrufen (Claude-Transkript), Start und Bestätigen des Kuratierens (Export), feste Wortlaute der Leitung im Transkript (unscharfer Vergleich); Erklären endet rechnerisch 2:50 nach dem Wortlaut; Nachher-Bogen und SUS aus den Nachbarn. Für 1.6: lokale Transkription und Vorschläge eines lokalen Modells als KI-Nutzung in der Auswertung; nichts verlässt den Rechner.

\section{Teilnehmende, Ablauf, Einwilligung, Pilot und Generalprobe}% 1,5 S.
\label{sec:teilnehmende}
% ZIEL: Wer teilnimmt, wie zugeteilt wird, wie eine Sitzung läuft, wie die Einwilligung zustande kam, was Pilot und erste Sitzungen am Ablauf geändert haben. Ablauftabelle (Minute, Teil) in den Text; Leitfaden, Ablaufvorlagen, Einwilligung in den Anhang.
%
% INHALT (geprüft)
% - Acht Studierende des Teamprojekts an teamable; Vorabeinschätzung: viel Arbeit mit KI-Agenten, wenig eigene Programmiererfahrung, der Vorfragebogen erhebt das (Plan). Zum Zeitpunkt der Sitzungen im zweiten Sprint, "weiß ich nicht" ist eine normale Antwort (Ablauf). Zwei Entwickler außerhalb des Teams als Vergleichsfälle, getrennt berichtet: zeigen, ob die Items nach Vorwissen trennen und wie jemand mit Erfahrung kuratiert; einmal D auf #74, einmal D auf E (Plan).
% - Zuteilung (Zut, festgelegt 12.10.): Pseudonym P01 bis P10 in der Reihenfolge des Rücklaufs der Einwilligung (Leit). Paar A = 74-C und E-D (P02, P03, P05, P07, P09), Paar B = 74-D und E-C (P01, P04, P06, P08, P10); P01 Entwickler 1, P02 Entwickler 2, P03 bis P10 abwechselnd. Reihenfolge der gewerteten Aufgaben je Paar dreimal die eine, zweimal die andere zuerst; #79 immer zuerst, immer unter D. Drei Kopien je Person.
% - Ablauf etwa 100 min (Ablauf, Stand 02.11.; Exposé: 60 min): 0 bis 10 Vorfragebogen (Aufnahme läuft, Bestätigung der Einwilligung auf der Aufnahme); 10 bis 13 Einführung Codebasis (App in drei Sätzen, Seiten Students, Team-Sets, Setup; im Explorer nur Ordner, src/hooks/ und students/ nicht aufklappen, die Dateinamen beantworten Items); 13 bis 28 Einführung Werkzeug und Aufwärmen #79 in der Reihenfolge der gewerteten Aufgaben: a Aufgabe lesen, b zwei Beispiel-Items mit Auflösung aus dem Kopf (ohne Provider-Muster, Prüf), c Oberfläche bis 5 min (der Proband klickt, die Leitung erklärt je Stufe; nach "Kontext bestätigen" Esc und /exit, der Lauf wird nicht abgewartet), d Anleitung zeigen; 28 bis 57 Aufgabe 1; 57 bis 86 Aufgabe 2; 86 bis 90 Reserve; 90 bis 100 SUS und Interview. Je gewertete Aufgabe sechs Schritte: Aufgabenblatt 2 min (nur der Bogen vorn, Browser und Editor hinten), Vorher-Bogen bis Minute 4 (Ansage bei 3), Lauf (Enter durch den Proband), Durchsehen genau 5, Erklären 2, Nachher-Bogen bis 5 (Ansage bei 4). C etwa 21 min, D bis 36 (Plan).
% - Umgebung (Plan E5, Zut, Prot 20./21.10.): Fernzugriff RustDesk (am 19.10. statt TeamViewer, Lizenzfrage), Ton über Discord, der Proband sieht den ganzen Bildschirm und nur das Fenster, das vorn ist; OBS; eigenes Chrome-Profil; VS-Code-Studienprofil ohne Claude-Code-Erweiterung, je Kopie ein Fenster mit Terminal; Dateinamen der Bögen ohne Aufgabe und Bedingung; Neo4j zwischen Probanden nicht neu starten; "Help improve Claude" im Studienkonto aus, Anthropic hält die Sitzungen 30 Tage (Leit).
% - Einwilligung (eval/studie/einwilligung.md): Aufzeichnung von Bildschirm und Ton, keine Kamera, Fernzugriff, Export, Agentenprotokoll, Bögen, Notizen; Übermittlung der Eingaben an Anthropic mit Standardvertragsklauseln (Art. 46 DSGVO; EU-US-Framework nicht bestätigt, Prüf), Löschung dort nach 30 Tagen, bei Sicherheitsprüfung bis zwei Jahre; Pseudonym, Zuordnungsliste getrennt, Erkennbarkeit in der kleinen Gruppe benannt, Vorfragebogen nur je Gruppe; Löschfristen, Rechte, Widerruf, Freiwilligkeit ohne Einfluss auf die Bewertung; wörtliche Zitate als eigener Punkt; Bitte um Verschwiegenheit bis zum letzten Termin. Weg: vorab unterschrieben (Scan, Foto) oder Mail "Einverstanden", Bestätigung zu Beginn auf der Aufnahme. HTWG: keine Hochschulvorlage, die Professur entscheidet (13.10.). Beschluss 19.10.: mündlich zuerst. Vorstellung im Team 29.10. (Ablauf, #79 als Beispiel, Beispiel-Items, Vorfragebogen A bis D, SUS; nicht gezeigt: gewertete Aufgaben und Teil E), Datenschutz mündlich, alle Anwesenden stimmten zu; Terminfindung über ein Google-Doc (Prot 29.10.). Zwei Rückläufe liegen (03.11., 05.11.). Drei Rückläufe (02.11., 05.11., 05.11.); verschickte Fassung in 6.5 WORTLAUT (07.11.).
% - Pilot PG 22.10. (eigene Kennung, Paar B; laut Gl die "Generalprobe mit Entwickler 1", damit die Studie bei P01 beginnt und keine Probandenkopie verbraucht wird; Prot 20.10.): Befunde: Aufgabenblatt 5,5 statt 2 min bei offenem Frontend, Vorher-Bogen 5 min (Grenze 4), Erklären am Aufwärmen entfiel, die Leitung klickte die Oberfläche durch (bei #74 brauchte der Proband dann die vollen 15 min), Einführung ohne konkrete Pfade, Streuzeichen in den Bögen, Pause wegen eines Notfalls (Mit PG, Prot 23.10.). Folgen: Grundwert gestrichen, beim Aufgabenblatt nur der Bogen vorn, Pfade im Ablauf, der Proband klickt selbst, laute Antworten. Ob die PG-Daten mitzählen: offen (Prot 20.10.). Kurzpilot PK: angelegt, entfällt, seit der ersten Sitzung gegenstandslos (Plan). Entwickler 2 (P02): Sitzungsordner 23.10., nicht gelaufen.
% - Änderungen 02.11. vor der ersten Sitzung (628b04f): je Schritt Vorn, Sagen, Handlung, Mitschrieb; Aufwärmen in der Reihenfolge Aufgabe, Items, Oberfläche; Chrome beim Erklären minimieren; Interviewfrage 4 in beide Richtungen mit Nachfrage; Teil A ans Ende der Bögen, TLX-Anker; Wortlaut der Items unverändert.
% - GEFAHRENER STAND (Mit; für 7.1 und das Nachziehen 30.11. bis 11.12.):
%   P03 03.11., Paar A (74-C, dann E-D): Vorfragebogen 5,5 min; Aufwärmen bis Minute 19 (Kuratieren 5); #74-C Vorher-Bogen etwa 5 min (über der Grenze), Lauf 2:33, Nachher 5; E-D Vorher 2:40, Kuratieren 15 (Maximum), Lauf 2, Nachher 6; Ende nach etwa 86 min; Internetabbruch bei der letzten Interviewantwort, wiederholt.
%   P04 05.11., Paar B (74-D, dann E-C): Zwischenablage defekt, Start-Aufgabe lieferte den Prompt nicht, der Proband kopierte den Aufgabentext selbst; 20 min Pause; Aufwärmen Kuratieren 7; #74-D Vorher 3 (alle Items "weiß ich nicht"), Kuratieren 14, Lauf 16 min, Nachher 8; E-C Vorher 2,3, Lauf 2:54, Nachher 6; Cap-Regler nach oben gezogen.
%   Folgerung für Kapitel 6: Zeitgrenzen der Bögen (4 und 5 min) werden nicht gehalten, Nachher-Bogen 5 bis 8 min, die Reserve reicht; je Person Dauer und Störungen in 7.1.
%
% VERWEISE: 6.2 (Technik je Bedingung), 7.1 (Abweichungen je Sitzung), 8.2, Anhang (Leitfaden, Ablauf, Einwilligung, Vorfragebogen, Anleitung für den Probanden).
% OFFEN: Entwickler-Sitzungen (P01, P02) und ob PG als Entwickler 1 zählt; zweite Gruppe: entfällt (29.10., nur dieses Teamprojekt); Einwilligung-Fassung: geklärt 07.11. (6.5 WORTLAUT); Reihenfolge der Pseudonyme (P03 lief vor P01 und P02, Plan-Satz "Reihenfolge des Rücklaufs" anpassen).

% GERÜST 6.5 (noch zu schreiben, 1,5 S.)
% Absatz 1: Teilnehmende (acht Studierende des Teamprojekts, zweiter Sprint; zwei Entwickler als Vergleichsfälle; Vorabeinschätzung; Vorfragebogen je Gruppe; Ausschlussregel PR-Autor; Bitte um Verschwiegenheit und D1 des Vorfragebogens als Kontrolle). Abweichung 100 statt 60 min in einem Satz mit Grund (Vorher-Nachher und Erklären kamen dazu; Aufwärmen).
% Absatz 2: Zuteilung (Pseudonym P01 bis P10; Paar A und B; Reihenfolge je Paar dreimal die eine, zweimal die andere; #79 immer zuerst unter D; drei Kopien). Hinweis: Pseudonyme wurden vor den Terminen vergeben, P03 lief vor P01 und P02; den Plan-Satz "Reihenfolge des Rücklaufs" nicht übernehmen, sondern "feste Zuteilung je Pseudonym, Termine in der Reihenfolge der Verfügbarkeit".
% Absatz 3: Ablauf als Tabelle tab:ablauf (Minute, Teil) und die sechs Schritte je gewerteter Aufgabe in einem Satz (stehen schon in 6.2). Fernzugriff, Ton, Aufnahme, was der Proband sieht, Eingabeweg der Bögen.
% Absatz 4: Einwilligung (Inhalt in zwei Sätzen, Weg: Vorstellung im Team 29.10. mit mündlicher Zustimmung aller, PDF per Mail, Rücklauf unterschrieben oder per Mail, Bestätigung auf der Aufnahme; keine Hochschulvorlage, Professur entscheidet; Fristen 11.12.2026 und 27.11.2026; Löschung 16.02.2027 und 30.09.2027). Übermittlung an Anthropic und Discord in einem Satz.
% Absatz 5: Pilot PG 22.10. und was er änderte (Grundwert gestrichen, nur der Bogen vorn beim Aufgabenblatt, Pfade in der Einführung, der Proband klickt selbst, laute Antworten) und die Änderungen vom 02.11. vor der ersten Sitzung. Ob PG mitzählt: Entscheidung nennen (Vorschlag: als Pilot getrennt, nicht in den acht, alter Ablauf). Stand der Sitzungen bis 07.11.: P03 und P04 gelaufen, Zeitgrenzen der Bögen nicht gehalten; Einzelheiten in 7.1.
% WORTLAUT 6.5, Zuteilung (kopien-zuteilung.ps1, festgelegt 12.10.): Paar A = 74-C und E-D: P02 (Entwickler 2), P03, P05, P07, P09. Paar B = 74-D und E-C: P01 (Entwickler 1), P04, P06, P08, P10. Reihenfolge der gewerteten Aufgaben: #74 zuerst bei P01, P03, P04, P07, P08 (und PG); E zuerst bei P02, P05, P06, P09, P10. PG = Paar B mit #74 zuerst (eigene Kennung für die Generalprobe, damit keine Probandenkopie verbraucht wird), PK = nur 74-C (Kurzpilot, entfallen).
% WORTLAUT 6.5, Tabelle tab:ablauf (ablauf-vorlage.md, Stand 02.11.): 0 bis 10 Vorfragebogen (Aufnahme läuft; auf die Aufnahme: "Die Einwilligung von <Pseudonym> liegt unterschrieben vor."); 10 bis 13 Einführung Codebasis (erst die App in drei Sätzen mit den Seiten Students, Team-Sets, Setup im Browser, dann der Explorer der #79-Kopie, nur Ordner: src/app/(app)/course/[courseId]/ mit je Seite einem Ordner (students, team-sets, setup, profiles) und page.tsx, daneben (components), (hooks), (table); layout.tsx dort als Rahmen um alle Kursseiten, src/app/layout.tsx als Rahmen um die App; src/app/(accounts)/accounts/ Anmeldung; src/components/ geteilte Bausteine, ui/ Grundbausteine; src/hooks/ Hooks zum Server (use-...), src/lib/ Hilfsfunktionen; src/hooks/ und students/ nicht aufklappen); 13 bis 28 Einführung Werkzeug und Aufwärmen #79 (a Aufgabe lesen im Terminal, "Lies die Aufgabe. Du musst sie nicht lösen, das macht gleich der Agent."; b zwei Beispiel-Items, 1 min, "'Weiß ich nicht' ist hier eine normale Antwort, ihr seid im zweiten Sprint. Die Fragen vorher zeigen nur, was du vorher schon weißt.", danach Auflösung aus dem Kopf; c Oberfläche bis 5 min, der Proband klickt, die Leitung sagt je Stufe ein, zwei Sätze: Entitäten "Das sind die Begriffe, die das Werkzeug aus der Aufgabe gezogen hat.", Seed-Suche "Die Startpunkte im Code; links der Quelltext zum angeklickten Knoten.", Traversal "Die Nachbarn der Startpunkte.", Ranking "Die Reihenfolge, in der der Agent die Stellen bekommt, oben das Wichtigste; was unter dem Cap steht, bekommt er nicht."; je Stufe eine Handlung (Begriff entfernen; Knoten anklicken, Haken weg und wieder hin; Nachbarn laden; Zeile ziehen, Vorschau, bestätigen); nach dem Bestätigen Esc und /exit, der Lauf wird nicht abgewartet; d Anleitung zeigen); 28 bis 57 Aufgabe 1; 57 bis 86 Aufgabe 2 (dazwischen Wechsel etwa 1 min: Chrome auf den anderen Tab, Fenster der zweiten Kopie nach vorn, nichts beenden); 86 bis 90 Reserve; 90 bis 100 SUS und Interview.
% WORTLAUT 6.5, Beispiel-Items #79 (beispiel-items.md): Stelle 1 "Der Rahmen, der um alle Seiten eines Kurses liegt (Kopf, Inhalt, Fuß)" (spielt eine Rolle; keine; weiß ich nicht; richtig: Rolle). Frage 1 "Der Text 'Teamable Analytics', den das Logo ersetzen soll: wo steht er im Code?" a Direkt in der oberen Navigation, vor den anderen Einträgen; b In einer Einstellungsdatei, aus der die Navigation ihn liest; c In einem eigenen Baustein, den die Navigation dort einbindet; d Weiß ich nicht; Sicherheit (richtig c, Logo).
% WORTLAUT 6.5, die sechs Schritte je gewerteter Aufgabe mit Ansagen (ablauf-aufgabe-vorlage.md, 02.11.): Vorbereitung (Chrome auf den Tab der Aufgabe, angemeldet, Einstiegsseite, dann nach hinten; Fenster der Kopie: Start-Aufgabe, /mcp grün, Strg+V, nicht Enter, nach hinten). 1 Aufgabenblatt 2 min, vorn nur der Vorher-Bogen mit dem Aufgabentext oben: "Lies die Aufgabe. Du musst sie nicht lösen, das macht gleich der Agent. Danach stelle ich dir ein paar Fragen dazu, bevor er loslegt. Wenn er fertig ist, schaust du dir sein Ergebnis an und erklärst es danach in zwei Minuten jemandem aus deinem Team. Es ist völlig in Ordnung, 'weiß ich nicht' zu sagen." Bei 2 min: "Weiter zu den Fragen." 2 Vorher-Bogen 2 min, Schluss bei 4: Teil B, der Proband sagt Antwort und Sicherheit, "Wenn du bei einer Frage länger überlegst, sag 'weiß ich nicht' und geh weiter.", bei 3 min "noch eine Minute", leere Items bleiben leer. 3 Lauf: "Jetzt schickst du dem Agenten die Aufgabe. Der Text steht schon im Terminal, du musst nur Enter drücken. Nicht Strg+Enter." Dann C- oder D-Ansage (6.2 WORTLAUT), Laufzeit in den Mitschrieb. 4 Durchsehen genau 5 min (Ansage in 6.4 WORTLAUT), Urteil und Grund in den Mitschrieb. 5 Erklären 2 min: "Ich schließe jetzt Terminal, Editor und Browser.", Kopie schließen, Chrome minimieren, Ablauf-Fenster vorn, Anweisung (6.4 WORTLAUT). 6 Nachher-Bogen 3 min, Schluss bei 5, bei 4 "noch eine Minute", nicht nachsehen lassen.
% WORTLAUT 6.5, Vorbereitung und Nachbereitung (ablauf-vorlage.md, Leitfaden 21.10.): 15 min vorher Neo4j Desktop und Ollama starten, eigenes VS Code und Privates schließen, sitzung-vor.ps1 -Pseudonym (bricht ab, wenn etwas fehlt; öffnet Dienste, Chrome mit beiden Anmeldeseiten, ein VS-Code-Fenster je Kopie und eines mit dem Sitzungsordner), in beiden Tabs anmelden, je Kopie Start-Aufgabe prüfen (2.1.251, Opus 5, /mcp verbunden, /exit), OBS mit Bildschirm, Mikrofon und Call-Ton, Probeaufnahme 10 s, RustDesk-Einmalpasswort notieren, dem Probanden den Screenshot zum Stummschalten schicken (rustdesk-stummschalten.jpeg: Bildschirm-Menü, Haken bei Stummschalten), Verbindungsart (direkt oder Relay) in den Mitschrieb, Dienste-Fenster minimiert, Ablauf vorn. Nachher: Aufnahme sichern; sitzung-nach.ps1 (beendet Fenster, setzt Backend-74-Datenbank zurück, verschiebt Exporte und Transkripte zur Person, Gate je gewerteter Kopie, zeigt Git-Befehle); Git je Kopie; "geliefert" markieren vor dem Öffnen der Bögen; Prüflisten; Auswertung nach studie\auswertung\README.md. Unterschriebene Einwilligungen liegen unter studie\einwilligung\, getrennt von Probandenordnern und Exporten. Dateinamen der Bögen im Sitzungsordner ohne Aufgabe und Bedingung (05-aufgabe1-vorher.md usw.).
% WORTLAUT 6.5, Mitschrieb-Vorlage (sitzung-anlegen.ps1): Kopf Datum, Kopien in Reihenfolge, Fernzugriff, Aufnahme gestartet; Vorfragebogen Dauer; Aufwärmen Kuratierzeit, Auffälligkeiten; je gewertete Aufgabe Vorher-Bogen Dauer, Kuratierzeit (nur D), Laufzeit, Besonderheiten im Lauf (Abbrüche, Rückfragen, Fehler), Urteil (übernehmen, nachbessern, verwerfen), Grund, Nachher-Bogen Dauer; SUS Dauer, Interview Stichworte je Frage 1 bis 4; Störungen (Verzögerung beim Fernzugriff; Fenster, Technik, Sonstiges).
% WORTLAUT 6.5, Einwilligung in der verschickten Fassung (PDF "Einwilligung_Studie_Peeters", für Anhang app:frageboegen; Namen und Anschrift beim Einbinden entscheiden): Titel "Einwilligung zur Teilnahme". Kopf: "Studie zur Zusammenarbeit mit einem KI-Coding-Agenten am teamable-Frontend. Masterarbeit von Joschka Peeters, Fakultät Informatik, HTWG Konstanz, betreut von Prof. Dr. Rainer Mueller. Sitzungsleitung: Joschka Peeters, jo741pee@htwg-konstanz.de. Verantwortlich für die Daten: Joschka Peeters, [Anschrift des Autors]." Dann: "Du bekommst diesen Text vor der Sitzung per Mail. Bitte lies ihn, frag bei Unklarheiten nach und schick ihn vor der Sitzung unterschrieben zurück (Scan oder Foto); eine Antwort per Mail mit 'Einverstanden' reicht auch, dann schreib bitte dazu, ob wörtliche Zitate erlaubt sind. Ohne deine Zustimmung findet die Sitzung nicht statt. Zu Beginn der Sitzung hält die Sitzungsleitung auf der Aufnahme fest, dass die Zustimmung vorliegt."
% Abschnitt "Worum es geht": "Du arbeitest etwa 100 Minuten mit einem KI-Coding-Agenten (Claude Code) an kleinen Aufgaben im teamable-Frontend. Bei einem Teil der Aufgaben siehst du im Browser, welche Stellen im Code der Agent bekommt, und kannst die Auswahl verändern. Zwischendurch beantwortest du kurze Fragen zum Code und erklärst mündlich, was geändert wurde. Es geht nicht um Geschwindigkeit und nicht darum, wie gut du programmierst. Untersucht wird, wie Menschen mit einem Coding-Agenten zusammenarbeiten und welche Rolle die Auswahl der Codestellen dabei spielt."
% Abschnitt "Was aufgezeichnet wird" (Liste): Der Bildschirm des Studienrechners während der ganzen Sitzung. Ton während der ganzen Sitzung: deine Stimme und die der Sitzungsleitung über den Call. Du wirst gebeten, laut zu denken. Die Tonaufnahme wird lokal auf dem Studienrechner verschriftet. Keine Kamera. Schalte die Kamera im Call bitte aus. Was du in der Oberfläche anklickst und auswählst, mit Zeiten. Das Protokoll des Agenten: deine Eingaben, seine Antworten, seine Änderungen am Code. Deine Antworten in den Fragebögen und das Gespräch am Ende. Notizen der Sitzungsleitung: Zeiten, ihr Urteil über das Ergebnis der Aufgabe, Störungen. Dann: "Du arbeitest per Fernzugriff (RustDesk) auf dem Rechner der Sitzungsleitung. Auf deinem eigenen Gerät wird nichts aufgezeichnet; dort laufen nur RustDesk und der Call. Die Fragebögen füllst du auf dem Studienrechner aus. Bildschirm und Eingaben laufen über RustDesk, der Ton über Discord. Was du in Claude Code eintippst, wird zur Verarbeitung an den Anbieter Anthropic in den USA geschickt; gib dort keine persönlichen Daten ein. Die Übermittlung stützt sich auf die Standardvertragsklauseln der EU-Kommission (Art. 46 DSGVO), die Anthropic in seiner Datenschutzerklärung nennt. Die Eingaben werden von Anthropic nicht zum Training von KI-Modellen verwendet; das ist im Studienkonto so eingestellt. Anthropic löscht die Sitzungen dann innerhalb von 30 Tagen; Eingaben, die eine automatische Sicherheitsprüfung auslöst, bis zu zwei Jahre. Für den Call ist Discord Netherlands BV verantwortlich; Discord überträgt Daten auch in die USA und stützt sich dabei auf die Standardvertragsklauseln und das EU-US Data Privacy Framework."
% Abschnitt "Pseudonym und Zugriff" (Liste): "Deine Daten bekommen ein Pseudonym (zum Beispiel P03). Die Liste, welcher Name zu welchem Pseudonym gehört, liegt getrennt von den Daten nur bei der Sitzungsleitung und wird nach Abgabe der Arbeit gelöscht." "In der Arbeit stehen Werte je Pseudonym, Beschreibungen des Vorgehens und einzelne wörtliche Sätze aus dem Gespräch, nie dein Name. Die Gruppe ist klein, und die Teilnehmenden kommen größtenteils aus demselben Teamprojekt. Wer das Team kennt, kann einzelne Personen daran vermutlich erkennen. Zitate werden so gewählt, dass sie keine Person erkennbar machen." "Die Daten liegen auf dem Studienrechner der Sitzungsleitung. Zugriff auf die Aufnahmen hat nur die Sitzungsleitung. Prof. Dr. Rainer Mueller sieht pseudonymisierte Daten, keine Aufnahmen und keine Namensliste. Die Angaben aus dem Vorfragebogen (Rolle, Erfahrung, Pull Requests) stehen in der Arbeit nur zusammengefasst je Gruppe, nie je Pseudonym." "Bildschirm- und Tonaufnahmen werden nach der Auswertung gelöscht, spätestens am 16.02.2027. Die übrigen pseudonymisierten Daten werden nach Abschluss des Prüfungsverfahrens gelöscht, spätestens am 30.09.2027." "Rechtsgrundlage ist deine Einwilligung (Art. 6 Abs. 1 Buchst. a DSGVO). Du hast das Recht auf Auskunft, Berichtigung, Löschung, Einschränkung der Verarbeitung und Datenübertragbarkeit und kannst dich beim Landesbeauftragten für den Datenschutz und die Informationsfreiheit Baden-Württemberg beschweren."
% Abschnitt "Freiwilligkeit und Rücktritt" (Liste): "Die Teilnahme ist freiwillig. Du kannst die Sitzung jederzeit ohne Angabe von Gründen abbrechen. Auf Wunsch werden die bis dahin aufgezeichneten Daten sofort gelöscht." "Du kannst deine Einwilligung jederzeit per Mail an jo741pee@htwg-konstanz.de widerrufen, ohne Begründung; dann werden Aufnahmen, Bögen und Protokolle gelöscht. Was bis zum Widerruf verarbeitet wurde, bleibt rechtmäßig. Bis zum 11.12.2026 gehen deine Daten dann auch nicht in die Arbeit ein. Danach bleiben Ergebnisse, die schon in der Arbeit stehen, ohne deinen Namen erhalten." "Ob du teilnimmst, abbrichst oder deine Daten zurückziehst, hat keinen Einfluss auf das Teamprojekt oder eine Bewertung. Wer das Teamprojekt bewertet, bekommt keine Aufnahmen, keine Namensliste und keine Angaben aus dem Vorfragebogen je Person." Abschnitt "Eine Bitte": "Bis zum letzten Sitzungstermin, dem 27.11.2026, bitte nichts über die Aufgaben und Fragen an andere Teilnehmende weitergeben."
% Abschnitt "Zustimmung" ("Bitte jeden Punkt ankreuzen, dem du zustimmst."): Ich habe den Text gelesen und konnte Fragen stellen. Ich bin mit der Aufnahme von Bildschirm und Ton einverstanden. Ich bin damit einverstanden, dass meine Eingaben, die Protokolle des Agenten und meine Antworten gespeichert und unter Pseudonym ausgewertet werden. Einzelne wörtliche Sätze aus dem Gespräch dürfen ohne Namen in der Arbeit stehen. Die Teilnahme hängt nicht an diesem Punkt. Ich nehme freiwillig teil und weiß, dass ich jederzeit abbrechen und meine Einwilligung widerrufen kann. Felder Name, Datum, Unterschrift, Pseudonym (trägt die Sitzungsleitung ein). Prüfung vom 11.10. hatte angemerkt: EU-US Data Privacy Framework für Anthropic nicht bestätigt; die verschickte Fassung nennt für Anthropic nur die Standardvertragsklauseln, für Discord beides.
% WORTLAUT 6.5, Vorstellung im Team 29.10. (Notizen der Sitzungsleitung, etwa 25 min): 0 bis 3 worum es geht; 3 bis 6 Vorfragebogen A bis D durchblättern, Teil E nur angekündigt; 6 bis 12 Aufwärmaufgabe #79 lesen, Oberfläche live in PG-79-D (alle vier Stufen, Esc nach dem Bestätigen), Ansage "In der Sitzung habt ihr dafür bis 15 Minuten und denkt laut. Eine der zwei gewerteten Aufgaben läuft ohne diese Oberfläche."; 12 bis 17 eine gewertete Aufgabe Schritt für Schritt mit den Beispiel-Items; 17 bis 19 SUS und Gespräch; 19 bis 23 Datenschutz mündlich nach den Punkten der Einwilligung; 23 bis 25 Frage in die Runde, Termine, Geräte. Nicht gezeigt: #74 und E, deren Bögen, Schlüssel, Prompts, Teil E des Vorfragebogens. Alle Anwesenden stimmten Teilnahme und Aufzeichnung mündlich zu. Geklärt: Call über Discord; nur dieses Teamprojekt, keine zweite Gruppe (mit dem Betreuer abgesprochen); RustDesk ohne Installation, Angebot des Studienkontos; Termine per Google-Doc (Werktage 30.10. bis 27.11., Zweistundenblöcke); Einwilligung in die Gruppe, Rücklauf privat als Foto oder PDF vor der Sitzung. Erstes Anschreiben an das Team (Entwurf 10.09.) hatte zehn Vorfragen zum Repo-Stand (Fork oder Branches, Commit 98ce1c8, Node-Version, tsc und lint, Beteiligung vor dem Teamprojekt, bekannte Stellen, KI-Werkzeuge, Zeitraum) und kündigte 100 min und zwei bis drei Aufgaben an.
% WORTLAUT 6.5, Pilot PG 22.10. (Mitschrieb): Aufnahme ab 14:14; Vorfragebogen geplant 7; Aufwärmen bis 5 geplant, Schritt 5 (Erklären) nicht durchgeführt, Leitung klickte durch; Aufgabe 1 PG-74-D: Aufgabenblatt 5,5 min (Frontend offen), Vorher-Bogen 5 min, Kuratierzeit nicht festgehalten (Bestätigung etwa 14:49, laut Export 15:00 min), kurze Pause wegen eines medizinischen Notfalls; Aufgabe 2 PG-E-C ohne Zeiten. Interview: 1 im Detail nicht immer klar, oberflächlich schon, git diff als Sicherheit; 2 Aufgabe 1 übergebenen Kontext angeschaut plus ein, zwei Stellen weiter (fetch), Aufgabe 2 unklar, teils MCP, teils eigene Suche; 3 nicht nach den Fragen gesucht, im Frontend überlegt, wo es sein könnte; 4 nichts gestört, man kann weiterklicken; nicht benutzen bei einfachen Aufgaben, gerade Frontend, das man hinterher anschauen kann; eher als automatischer Ablauf; wenn der Agent falsch bezogen hat, selbst Kontext rausstreichen; wichtig, dass man Kontext wegnehmen kann. Streuzeichen in Vorfragebogen und beiden Bögen von Aufgabe 1 (Kreuze neben dem Kästchen, Zeichen im Text), mehrere Antworten nicht eindeutig; Ursache CRLF in den Vorlagen gegen die Positionsrechnung der Markdown-Vorschau, seit 23.10. werden die Dateien mit LF geschrieben.
% WORTLAUT 6.5, P01 (Entwickler 1): 21./22.10. nur Durchspielen durch den Autor (VS-Code-Fenster je Kopie, Sitzungsordner je Person, unrunde Stellen; Sitzungs-JSONs 21.10. 11:29, 11:48, 15:52 und 22.10. 09:48), Mitschrieb leer, keine Sitzung mit Probanden. P02 (Entwickler 2): Sitzungsordner 23.10., nicht gelaufen. Beide Termine offen.

\section{Auswertung}% 1,5 S.
% Inhalt: beschreibend und explorativ angelegt, acht Differenzen, Vorzeichentest ergänzend, zwei Schichten (geliefert, nicht geliefert), Kontrollitems gegen den Vortest als Suchauftrag, Entwickler getrennt.
\label{sec:auswertung}
% ZIEL: Der Leser weiß vor Kapitel 7, welche Zahl je Frage berichtet wird, wie sie entsteht und was sie nicht leisten kann. Keine Ergebnisse. Eine Tabelle: Frage, Kennzahl, Vergleich, Darstellung in Kapitel 7.
%
% INHALT (geprüft)
% - Beschreibend, die Werte aller acht Personen einzeln; in Kapitel 7 Tabellen statt Prosa, je Person eine Zeile (Plan, Gl).
% - Hauptkennzahl Frage 2: mittlerer Zuwachs (nachher minus vorher) je nicht geliefertem Item, Anzahl dieser Items daneben (3 bis 5 je Sitzung, nach Lösungsweg des Agenten; Prüf), in D minus in C je Person, je Aufgabe zentriert (Mittel aller acht auf dieser Aufgabe abgezogen), acht Differenzen. Vorzeichentest ergänzend (8:0 ergibt 0,008, 7:1 ergibt 0,07). Berichtet wird die Richtung, kein Nachweis (Plan, Mess).
% - Daneben: Vorher-Werte D gegen C je Aufgabe (Gruppen vor der Behandlung); Zuwachs D gegen C je Aufgabe, 4 gegen 4; Nachher-Werte der nicht vorgefragten Items: liegt D nur bei den vorgefragten vorn, hat der Vortest gewirkt; Kontrollitems E5 und #74-6 (wo sie nicht geliefert wurden) (Plan).
% - Zwei Schichten: je Sitzung vor dem Lesen der Bögen aus Transkript und Diff markiert, welche Items die Zusammenfassung beantwortet; berichtet C gegen D bei gelieferten und bei nicht gelieferten Items (Plan, Leit).
% - Frage 1: kuratiert minus neu gerechneter automatischer Kontext je Person, auf Symbolebene und als Deckung, dazu umgekehrter Anteil und Größe; Eingriffe je Stufe, Dauer, Abstand der Reihenfolge; je Aufgabe getrennt wegen des Deckeneffekts bei #74 (Plan).
% - Verbindung: je Item und Person der Status des Knotens (in D: nicht gezeigt, gezeigt, angesehen, behalten, abgewählt; in beiden: vom Agenten gelesen, im Diff, in der Zusammenfassung) gegen richtig vorher und nachher. Muster, kein Test (Plan).
% - Kalibrierung: Anteil sicher (4 oder 5) und falsch, je Bedingung, vorher und nachher. Empfindlichkeit daneben: dieselben Zuwächse mit falsch minus ein Drittel und als sichere Trefferquote; Anteil "weiß ich nicht" und Sicherheit 1 oder 2 je Bedingung. Hauptkennzahl bleibt die Wertung aus E3 (Plan).
% - Erklären: übernommen und eigenständig je Person und Bedingung, ohne Grundwert (Plan).
% - Nebenfragen: Prüfliste C gegen D je Aufgabe (FF4); Kuratierzeit, TLX-Rohwerte, SUS (FF6); Interview thematisch mit Zitaten, den Ereignissen im Export zugeordnet (Plan). E3 aus views, Eingriffen, lautem Denken und Interviewfrage 3.
% - Vorfragebogen nur je Gruppe, nie je Pseudonym: der Betreuer bewertet das Teamprojekt, Rolle und PR-Zahl wären bei acht Personen zuordenbar (Plan, Prüf).
% - Die zwei Entwickler getrennt. "Kein erkennbarer Unterschied" ist ein vollwertiges Ergebnis: Kuratieren bringt dann kein Verständnis über eine gute Zusammenfassung hinaus (Plan, Mess). Was das Design nicht zeigen kann: einen Nachweis, die Trennung von Sehen und Wählen, Verständnis auf Zeilenebene (Mess).
% - Hilfsmittel: kontext-gegen-referenz.js (Frage 1), aufnahme-auswerten.py und sprecher-zuordnen.py (Transkript, Zeitleiste, Vorschläge für Erklären und Interview; Ausw), kein Auslese-Skript für die Bögen (07.11. geprüft), Übertragung von Hand. Welche Schritte ein lokales Modell vorschlägt und dass die Wertung bei der Leitung bleibt, hier oder in 1.6 sagen.
%
% VERWEISE: 5.1 (Vorzeichentest), 6.4 (Wertung je Item, Schichten), 7.2 bis 7.5, 8.1, 8.2.
% OFFEN: Umgang mit PG und den Entwicklern in den Tabellen (eigene Zeilen, nicht im Mittel); n, falls eine zweite Gruppe dazukommt; Schwelle für "angesehen" (views ab 300 ms) festlegen, bevor die Exporte ausgewertet werden.
% GERÜST 6.6 (noch zu schreiben, 1,5 S.)
% Absatz 1: Grundsatz: beschreibend, Werte aller Personen einzeln, in Kapitel 7 Tabellen je Person; Richtung statt Nachweis; die zwei Entwickler und PG getrennt (eigene Zeilen, nicht im Mittel); "kein erkennbarer Unterschied" als vollwertiges Ergebnis.
% Absatz 2: Hauptkennzahl Frage 2 mit der Rechnung (Zuwachs je nicht geliefertem Item, D minus C je Person, je Aufgabe zentriert, acht Differenzen; Vorzeichentest ergänzend mit 8:0 p=0,008 und 7:1 p=0,07, Verweis 5.1). Daneben Vorher-Werte D gegen C, Zuwachs je Aufgabe 4 gegen 4, nicht vorgefragte Items, Kontrollitems, Kalibrierung mit Empfindlichkeit, Erklären.
% Absatz 3: Frage 1 (kuratiert minus neu gerechnet, Symbolebene und Deckung, Eingriffe je Stufe, Dauer, Abstand der Reihenfolge; je Aufgabe getrennt) und die Verbindung je Item und Person (Status des Knotens: nicht gezeigt, gezeigt, angesehen ab 300 ms, behalten, abgewählt; in beiden Bedingungen: gelesen, im Diff, in der Zusammenfassung).
% Absatz 4: Nebenfragen (Prüfliste C gegen D, Kuratierzeit, Raw TLX, SUS nach Brooke, Interview thematisch mit Zitaten, den Ereignissen im Export zugeordnet; E3 aus views, Eingriffen, lautem Denken, Interviewfrage 3). Vorfragebogen je Gruppe, Grund.
% Absatz 5: Hilfsmittel und KI-Kennzeichnung (kontext-gegen-referenz.js; Transkription und Vorschläge lokal; Bögen von Hand ausgelesen, kein Skript); was das Design nicht zeigen kann (Nachweis, Trennung von Sehen und Wählen, Zeilenebene).
% Tabelle tab:auswertung (Frage, Kennzahl, Vergleich, Darstellung in Kapitel 7): Frage 2 / Zuwachs je nicht geliefertem Item / D minus C je Person, zentriert / Tabelle je Person, Vorzeichen; Frage 2 / Erklären übernommen und eigenständig / je Bedingung / Tabelle je Person; Frage 2 / Kalibrierung sicher und falsch / je Bedingung, vorher und nachher / Tabelle; Frage 1 / Deckung und Symbolebene kuratiert gegen neu gerechnet / je Person, je Aufgabe / Tabelle je Person; Verbindung / Knotenstatus gegen richtig / je Item und Person / Matrix, Muster; FF4 / Prüfliste bestanden, Überschneidung / C gegen D je Aufgabe / Tabelle; FF6 / Kuratierzeit, Raw TLX, SUS / je Bedingung bzw. je Person / Tabelle; FF6 / Interview / thematisch / Zitate; E3 / angesehene Knoten, Eingriffe, Begründung geöffnet / je Person / beschreibend.
% WORTLAUT 6.6, Vorfragebogen-Schlüssel: E1 a, E2 c, E3 b; Summe 0 bis 3 als Niveau Codelesen; daneben A1 bis A3, B1 bis B3, C3 beschreibend; Gruppen Entwickler und Studierende getrennt, innerhalb der Studierenden nur zur Einordnung; nie je Pseudonym (Einwilligung; der Betreuer bewertet das Teamprojekt).
% WORTLAUT 6.6, Datenlage je Sitzung (was auf dem Studienrechner liegt, für Kapitel 7): Sitzungsordner <Pseudonym>-sitzung (00-ablauf, 01-mitschrieb, 02-vorfragebogen, 03-anleitung, 04-aufwaermen-items, 05 bis 08 Bögen in Sitzungsreihenfolge, 09-sus, seit P03 10-/11-pruefliste-<Aufgabe>-<Bedingung> mit Screenshots und 12-/13-geliefert-<Aufgabe>-<Bedingung>); <Pseudonym>-export (JSON der Studienseite je D-Aufgabe); <Pseudonym>-transkripte (Claude-Code-Transkripte); <Pseudonym>-gate-<Aufgabe>-<Bedingung>.txt; Kopien als Branches im Bare-Repo; OBS-Aufnahme im Videoordner; auswertung\<Pseudonym>\ mit Transkript, Zeitleiste, Vorschlägen. Stand 07.11.: vollständig für P03 und P04, PG ohne Prüflisten und geliefert-Dateien.
